% !TeX root=../main.tex
\chapter{پیاده‌سازی عملی، ساخت نمونه و آزمون‌های واقعی}

\section{معماری سامانه عملی اسکوتر خودران}
\label{sec:real_scooter_architecture}

\subsection{نمای کلی معماری سه‌سطحی سامانه}

در این فصل، روند پیاده‌سازی عملی سامانه پیشنهادی، آماده‌سازی نمونه واقعی اسکوتر، ارتباط میان واحدهای سخت‌افزاری و نرم‌افزاری، و نحوه اجرای آزمون‌های واقعی تشریح می‌شود. هدف اصلی این فصل، تبیین چگونگی انتقال ساختار طراحی‌شده در محیط شبیه‌سازی به یک سامانه فیزیکی قابل آزمون است؛ به‌گونه‌ای که داده‌های حسگرهای واقعی به قالب مشاهده مورد استفاده عامل یادگیری تقویتی تبدیل شده و فرمان تولیدشده توسط مدل آموزش‌دیده به فرمان قابل اعمال به موتورهای اسکوتر تبدیل شود.

معماری سامانه عملی از سه سطح اصلی تشکیل شده است. در سطح پایین، آردوینو مگا وظیفه ارتباط مستقیم با انکودرهای چهار موتور، راه‌اندازهای موتور و پنج حسگر فراصوت را بر عهده دارد. این واحد داده‌های مربوط به وضعیت چرخ‌ها و فاصله‌های نزدیک‌برد حسگرهای فراصوت را دریافت کرده و فرمان نهایی حرکت را به راه‌اندازهای موتور منتقل می‌کند. در سطح میانی، رزبری پای ۵ به‌عنوان گره واسط حسگری و ارتباطی عمل می‌کند. این واحد داده‌های تصویری دوربین استریو، اطلاعات واحد اندازه‌گیری اینرسی و داده‌های دریافتی از آردوینو را جمع‌آوری کرده و آن‌ها را برای پردازش به لپ‌تاپ ارسال می‌کند. در سطح بالا، لپ‌تاپ وظیفه پردازش لحظه‌ای داده‌ها، تشکیل بردار مشاهده، اجرای مدل تصمیم‌گیر آموزش‌دیده و تولید فرمان حرکتی را بر عهده دارد. در این ساختار، مدل یادگیری تقویتی در مرحله عملی صرفاً در حالت استنتاج اجرا شده و فرایند آموزش جدیدی روی سامانه واقعی انجام نمی‌شود.

تفکیک معماری به سه سطح یادشده، با هدف جداسازی وظایف کنترل پایین‌سطح، جمع‌آوری و انتقال داده‌های حسگری، و پردازش و تصمیم‌گیری سطح‌بالا انجام شده است. به این ترتیب، وظایفی مانند خواندن انکودرها، دریافت داده‌های حسگرهای فراصوت و اعمال فرمان به موتورها در لایه پایین‌سطح باقی می‌مانند، درحالی‌که پردازش‌های سنگین‌تر مانند تخمین عمق، تشخیص هدف و موانع، تشکیل مشاهده و اجرای سیاست تصمیم‌گیر در لپ‌تاپ انجام می‌شوند. این ساختار علاوه بر کاهش بار پردازشی واحدهای تعبیه‌شده، امکان توسعه و آزمون جداگانه هر بخش از سامانه را نیز فراهم می‌کند.

شکل \ref{fig:practical_system_data_command_flow} نمای کلی جریان داده و فرمان در معماری سامانه عملی را نشان می‌دهد. همان‌طور که در این شکل مشاهده می‌شود، مسیرهای داده با فلش آبی و مسیرهای فرمان با فلش نارنجی مشخص شده‌اند. داده‌های دوربین استریو و واحد اندازه‌گیری اینرسی مستقیماً به رزبری پای ۵ وارد می‌شوند، درحالی‌که داده‌های انکودرهای چرخ و حسگرهای فراصوت ابتدا توسط آردوینو دریافت شده و سپس از طریق ارتباط \lr{USB} به رزبری پای منتقل می‌شوند. در ادامه، داده‌های تجمیع‌شده از طریق شبکه به لپ‌تاپ ارسال شده و پس از پردازش، فرمان حرکتی تولیدشده از مسیر لپ‌تاپ، رزبری پای، آردوینو مگا و راه‌اندازهای موتور به موتورهای اسکوتر اعمال می‌گردد.

\begin{figure}[ht]
	\centerline{\includegraphics[width=0.75\linewidth]{robot_sensor_and_control_system_diagram.png}}
	\caption{جریان کلی داده و فرمان در سامانه عملی اسکوتر خودران}
	\label{fig:practical_system_data_command_flow}
\end{figure}


\subsection{ساختار پردازشی و کنترلی سامانه}
\label{subsec:processing_control_structure}

در پیاده‌سازی عملی، مسیرهای ارتباطی سامانه به دو بخش اصلی تقسیم شدند: مسیر انتقال داده‌های تصویری و مسیر انتقال داده‌های عددی و فرمان‌های حرکتی. داده‌های تصویری دوربین استریو به دلیل حجم بالا، به‌صورت جریان‌های جداگانه چپ و راست و با استفاده از \lr{GStreamer} روی شبکه محلی به لپ‌تاپ ارسال می‌شوند. در مقابل، داده‌های عددی مانند اودومتری، مقادیر حسگرهای فراصوت و سرعت موتورها از طریق بسته‌های \lr{UDP} میان رزبری پای و لپ‌تاپ منتقل می‌گردند.

در سطح میانی، رزبری پای علاوه بر انتقال داده‌ها، بخشی از پردازش حرکتی را نیز انجام می‌دهد. داده‌های انکودر و حسگرهای فراصوت از آردوینو دریافت می‌شوند و داده \lr{yaw} از واحد اندازه‌گیری اینرسی خوانده می‌شود. سپس اودومتری نسبی اسکوتر با ترکیب داده‌های انکودر و \lr{IMU} محاسبه شده و همراه با مقادیر حسگرهای فراصوت و سرعت موتورها برای لپ‌تاپ ارسال می‌شود. در سمت مقابل، فرمان تولیدشده توسط مدل تصمیم‌گیر در لپ‌تاپ ابتدا به سرعت هدف چهار موتور تبدیل شده و سپس از طریق \lr{UDP} به رزبری پای فرستاده می‌شود. رزبری پای این فرمان را به قالب سریال قابل فهم برای آردوینو تبدیل کرده و از طریق \lr{USB/Serial} برای اجرا ارسال می‌کند.

مسیر فرمان از دید اجرایی به این صورت است که مدل \lr{LSTM-SAC} روی لپ‌تاپ فرمان سرعت خطی و زاویه‌ای را تولید می‌کند. سپس گره واسط سخت‌افزاری لپ‌تاپ این فرمان را به مقادیر \lr{RPM} چهار موتور تبدیل می‌نماید. این مقادیر از طریق \lr{UDP} به رزبری پای منتقل شده و در نهایت با قالب \serialcmd  به آردوینو ارسال می‌شوند. آردوینو نیز با استفاده از بازخورد انکودرها، کنترل سرعت پایین‌سطح موتورهای چهارگانه را انجام می‌دهد.

جدول \ref{tab:real_system_data_command_paths} مسیرهای اصلی تبادل داده و فرمان در سامانه عملی را خلاصه می‌کند.

\begin{table}[htbp]
	\centering
	\caption{مسیرهای اصلی تبادل داده و فرمان در سامانه عملی}
	\label{tab:real_system_data_command_paths}
	\begin{tabular}{|p{0.2\textwidth}|p{0.58\textwidth}|}
		\hline
		\textbf{مسیر ارتباطی} &
		\textbf{محتوای منتقل‌شده و بستر ارتباطی}\\
		\hline
		
		دوربین استریو به لپ‌تاپ &
		انتقال جریان‌های تصویری چپ و راست با \lr{GStreamer} روی \lr{LAN}. \\
		\hline
		
		آردوینو به رزبری پای &
		ارسال داده انکودرهای چهار موتور و مقادیر حسگرهای فراصوت از طریق \lr{USB/Serial}. \\
		\hline
		
		رزبری پای به لپ‌تاپ &
		ارسال تله‌متری شامل اودومتری نسبی، داده‌های فراصوت و سرعت موتورها با \lr{UDP} روی \lr{LAN}. \\
		\hline
		
		لپ‌تاپ به رزبری پای &
		ارسال فرمان سرعت هدف چهار موتور پس از تبدیل \lr{cmd\_vel} به \lr{RPM}، با \lr{UDP} روی \lr{LAN}. \\
		\hline
		
		رزبری پای به آردوینو &
		ارسال فرمان نهایی چهارموتوره در قالب \serialcmd{} از طریق \lr{USB/Serial}. \\
		\hline
	\end{tabular}
	
\end{table}




\subsection{لایه حسگری و انتقال داده‌ها}
\label{subsec:sensing_and_data_transfer}

لایه حسگری سامانه عملی با هدف فراهم‌کردن اطلاعات لازم برای ناوبری هدف‌محور، تشخیص موانع، پایش فاصله‌های نزدیک‌برد و تخمین وضعیت حرکتی اسکوتر طراحی شد. در این سامانه، دوربین استریو، حسگرهای فراصوت، انکودرهای موتور و واحد اندازه‌گیری اینرسی به‌عنوان منابع اصلی داده مورد استفاده قرار گرفتند. خروجی این حسگرها پس از انتقال به واحدهای پردازشی، در تشکیل مشاهده عامل، تخمین اودومتری، لایه ایمنی و ثبت داده‌های آزمون به کار گرفته می‌شوند.

برای دریافت اطلاعات دیداری محیط، از دوربین استریو مبتنی بر دو حسگر \lr{IMX219-83} استفاده شد. هر دو حسگر از طریق کابل \lr{CSI} به رزبری پای متصل هستند و تصاویر چپ و راست ابتدا روی رزبری دریافت می‌شوند. تصاویر با وضوح ($1640 \times 1232$) و با استخراج صفحه روشنایی، به‌صورت خاکستری در قالب \lr{mono8} آماده می‌شوند. سپس دو جریان تصویری مستقل با نرخ تقریبی ۱۵ فریم بر ثانیه، از طریق \lr{GStreamer} و بستر \lr{LAN} به لپ‌تاپ منتقل می‌گردند. این داده‌ها در لپ‌تاپ برای تصحیح هندسی، تولید نقشه اختلاف مکانی، تخمین عمق، تشخیص هدف و تشخیص موانع استفاده می‌شوند.

حسگرهای فراصوت \lr{SRF05} برای پوشش نواحی نزدیک اطراف اسکوتر به کار گرفته شدند. پنج حسگر فراصوت در موقعیت‌های جلو-چپ، جلو-راست، چپ، راست و عقب نصب شده‌اند و همگی به آردوینو متصل هستند. آردوینو فاصله اندازه‌گیری‌شده توسط این حسگرها را دریافت کرده و از طریق ارتباط \lr{USB/Serial} به رزبری پای منتقل می‌کند. سپس این داده‌ها در قالب تله‌متری از طریق \lr{UDP} به لپ‌تاپ ارسال می‌شوند. داده‌های فراصوت هم در فضای مشاهده عامل و هم در منطق ایمنی نزدیک‌برد مورد استفاده قرار می‌گیرند.

انکودرهای چهار موتور برای پایش سرعت چرخ‌ها و پشتیبانی از کنترل پایین‌سطح استفاده شدند. هر موتور دارای انکودر مستقل است و آردوینو با خواندن پالس‌های انکودر، سرعت چرخ‌ها را محاسبه می‌کند. این داده‌ها علاوه بر استفاده در کنترل سرعت موتورهای چهارگانه، برای تخمین سرعت حرکتی اسکوتر نیز در اختیار رزبری پای قرار می‌گیرند.

با توجه به اینکه اسکوتر دارای ساختار حرکتی رانش لغزشی یا \lr{Skid-steering} است، استفاده از انکودرها به‌تنهایی برای تخمین جهت‌گیری در همه شرایط کافی نیست. در حرکت‌های همراه با چرخش یا لغزش، اختلاف میان چرخش واقعی بدنه و تخمین مبتنی بر اختلاف سرعت چرخ‌ها افزایش می‌یابد. به همین دلیل، از واحد اندازه‌گیری اینرسی مدل \lr{TSS-USB-S v2.0} برای دریافت زاویه \lr{yaw} استفاده شد. این حسگر از طریق \lr{USB} به رزبری پای متصل است و در ابتدای هر آزمون، مقدار اولیه \lr{yaw} آن به‌عنوان مرجع صفر در نظر گرفته می‌شود.

در پیاده‌سازی عملی، رزبری پای داده‌های انکودر و زاویه \lr{yaw} را با یکدیگر ترکیب می‌کند و اودومتری نسبی اسکوتر را می‌سازد. بنابراین، داده خام \lr{IMU} به‌صورت مستقیم وارد مشاهده عامل نمی‌شود؛ بلکه خروجی ترکیب‌شده شامل موقعیت نسبی، زاویه جهت‌گیری، سرعت خطی و سرعت زاویه‌ای از طریق تله‌متری \lr{UDP} به لپ‌تاپ ارسال می‌شود. این داده‌ها در تشکیل مشاهده و همچنین مدیریت حافظه موانع نسبت به حرکت اسکوتر استفاده می‌شوند.

مسیر انتقال داده‌ها به‌گونه‌ای طراحی شد که جریان‌های تصویری و داده‌های عددی از یکدیگر جدا باشند. تصاویر استریو به دلیل حجم بالا با \lr{GStreamer} منتقل می‌شوند، در حالی که داده‌های عددی مانند اودومتری، مقادیر حسگرهای فراصوت و سرعت موتورها از طریق بسته‌های \lr{UDP} روی شبکه محلی ارسال می‌گردند. این تفکیک باعث کاهش تداخل میان انتقال تصویر و داده‌های کنترلی می‌شود.

جدول \ref{tab:real_sensor_data_paths} خلاصه منابع حسگری و مسیر انتقال داده آن‌ها را نشان می‌دهد.

\begin{table}[htbp]
	\centering
	\caption{خلاصه منابع حسگری و مسیر انتقال داده در سامانه عملی}
	\label{tab:real_sensor_data_paths}
	\begin{tabular}{|p{0.24\textwidth}|p{0.31\textwidth}|p{0.35\textwidth}|}
		\hline
		\textbf{منبع داده} &
		\textbf{داده اصلی} &
		\textbf{مسیر انتقال و کاربرد} \\
		\hline
		دوربین استریو &
		تصاویر خاکستری چپ و راست &
		اتصال به رزبری پای از طریق \lr{CSI} و انتقال به لپ‌تاپ با \lr{GStreamer} برای تخمین عمق، تشخیص هدف و تشخیص موانع. \\
		\hline
		
		حسگرهای فراصوت &
		فاصله نزدیک‌برد در پنج جهت &
		اتصال به آردوینو و انتقال از مسیر رزبری پای به لپ‌تاپ برای مشاهده عامل و لایه ایمنی. \\
		\hline
		
		انکودرهای موتور &
		سرعت چرخ‌های چهارگانه &
		خواندن توسط آردوینو و استفاده در کنترل سرعت موتور و تخمین حرکت اسکوتر. \\
		\hline
		
		واحد اندازه‌گیری اینرسی &
		زاویه نسبی \lr{yaw} &
		اتصال به رزبری پای از طریق \lr{USB} و ترکیب با داده انکودر برای تولید اودومتری نسبی. \\
		\hline
		
	\end{tabular}
	
	
\end{table}


\subsection{جریان فرمان در سامانه عملی}
\label{subsec:real_command_flow}

جریان فرمان در سامانه عملی از لایه تصمیم‌گیری سطح‌بالا در لپ‌تاپ آغاز شده و پس از عبور از رزبری، آردوینو و راه‌اندازهای موتور، به حرکت چرخ‌های اسکوتر تبدیل می‌شود. مدل تصمیم‌گیر \lr{LSTM-SAC} پس از دریافت مشاهده، یک کنش دوبعدی نرمال‌شده تولید می‌کند که مطابق رابطه \ref{eq:action_scaling_final} به سرعت خطی و زاویه‌ای مطلوب اسکوتر تبدیل می‌شود. بنابراین، خروجی مستقیم مدل فرمان \lr{PWM} یا فرمان مستقل موتورها نیست، بلکه یک فرمان حرکتی سطح‌بالا در قالب (v) و ($\omega$) است.

فرمان (v) و ($\omega$) در لپ‌تاپ و در قالب پیام سرعت منتشر می‌شود. سپس گره واسط سخت‌افزاری لپ‌تاپ این فرمان را بر اساس مدل حرکتی اسکوتر به سرعت هدف چهار موتور تبدیل می‌کند. این مقادیر از طریق ارتباط \lr{UDP} روی شبکه محلی به رزبری پای ارسال می‌شوند. رزبری پس از دریافت فرمان، آن را به قالب سریال قابل فهم برای آردوینو تبدیل کرده و از طریق ارتباط \lr{USB/Serial} به آردوینو می‌فرستد.

آردوینو پس از دریافت فرمان، مقدار هدف هر موتور را استخراج کرده و کنترل سرعت پایین‌سطح موتورهای چهارگانه را با استفاده از بازخورد انکودرها انجام می‌دهد. در این مرحله، فرمان هدف موتور به سیگنال مناسب برای راه‌اندازهای موتور تبدیل می‌شود. بنابراین، آردوینو صرفاً یک ارسال‌کننده مستقیم فرمان نیست، بلکه وظیفه بستن حلقه کنترل سرعت موتورها را نیز بر عهده دارد.

در مسیر فرمان، محدودیت‌های ایمنی نیز اعمال می‌شوند. در شرایط عادی، فرمان تولیدشده توسط عامل پس از تبدیل و انتقال اجرا می‌شود؛ اما در صورت تشخیص وضعیت خطر توسط لایه ایمنی، فرمان نهایی می‌تواند محدود یا متوقف شود. آستانه‌ها و نواحی خطر مطابق منطق سپر ایمنی معرفی‌شده در بخش \ref{subsec:safety_shield_levels_final} و جدول \ref{tab:shield_thresholds_final} در نظر گرفته شده‌اند و جزئیات اجرای آن در بخش \ref{subsec:real_safety_logic_execution} توضیح داده می‌شود.

همچنین برای جلوگیری از باقی‌ماندن آخرین فرمان در شرایط خطا، در مسیر ارتباطی فرمان سازوکار توقف خودکار در نظر گرفته شد. در صورتی که گره واسط لپ‌تاپ فرمان جدیدی از مدل دریافت نکند، یا رزبری پای بسته فرمان تازه‌ای از لپ‌تاپ دریافت نکند، فرمان صفر برای موتورها ارسال می‌شود. به این ترتیب، قطع ارتباط، توقف نرم‌افزار یا تأخیر غیرعادی باعث ادامه نامحدود آخرین فرمان حرکتی نمی‌شود.



\subsection{ارتباط معماری واقعی با معماری شبیه‌سازی}
\label{subsec:sim_real_architecture_relation}

معماری سامانه عملی بر پایه همان منطق کلی حلقه ادراک، تصمیم‌گیری و کنترل طراحی شد که در فصل چهارم برای محیط شبیه‌سازی معرفی شده بود. در هر دو بستر، ابتدا داده‌های محیطی دریافت و پردازش می‌شوند، سپس مشاهده عامل ساخته می‌شود، عامل تصمیم‌گیر فرمان حرکتی را تولید می‌کند و در نهایت فرمان پس از اعمال محدودیت‌های حرکتی و ایمنی به سامانه حرکتی منتقل می‌شود. حفظ این ساختار برای انتقال مدل آموزش‌دیده از شبیه‌سازی به سامانه واقعی ضروری است، زیرا عامل یادگیری تقویتی با قالب مشخصی از مشاهده و عمل آموزش دیده و در اجرای عملی نیز باید داده‌هایی با مفهوم، ترتیب، مقیاس و محدوده سازگار دریافت کند.

هدف از پیاده‌سازی عملی، تغییر منطق اصلی سامانه شبیه‌سازی‌شده نبود؛ بلکه تلاش شد همان زنجیره ادراک، تشکیل مشاهده، تصمیم‌گیری، تولید فرمان و اعمال سپر ایمنی در بستر واقعی بازسازی شود. تفاوت اصلی دو بستر در منبع داده‌ها و شرایط اجرای آن‌ها است. در شبیه‌سازی، داده‌ها از محیط مجازی و حسگرهای مدل‌شده به دست می‌آیند، در حالی که در سامانه عملی، داده‌ها از حسگرهای واقعی با نویز، تأخیر، محدودیت میدان دید، خطای کالیبراسیون و محدودیت نرخ پردازش دریافت می‌شوند.

یکی از تفاوت‌های مهم، روش تشخیص هدف است. در محیط شبیه‌سازی، هدف بر اساس ویژگی رنگی آن از تصویر جدا می‌شد؛ اما در سامانه عملی، تصاویر استریو به‌صورت خاکستری منتقل می‌شوند و اطلاعات رنگی برای تشخیص هدف در دسترس نیست. به همین دلیل، در اجرای واقعی از نشانگر \lr{AprilTag} استفاده شد. پس از تشخیص نشانگر، فاصله و زاویه نسبی هدف از ترکیب موقعیت تگ در تصویر و نقشه اختلاف مکانی استخراج می‌شود. سپس مشابه روشی که برای موانع استفاده می‌شود، مختصات نسبی هدف در زمان‌های متوالی با فیلتر کالمن پالایش شده و سرعت نسبی هدف محاسبه می‌گردد. بنابراین، اگرچه روش استخراج هدف از تصویر در سامانه عملی نسبت به شبیه‌سازی تغییر کرده است، خروجی نهایی همچنان شامل فاصله، زاویه و مؤلفه‌های سرعت نسبی هدف در قالب مورد نیاز عامل است.

تفاوت مهم دیگر به تخمین وضعیت حرکتی اسکوتر مربوط است. در محیط شبیه‌سازی، موقعیت و جهت‌گیری نسبی ربات از خروجی اودومتری و وضعیت منتشرشده توسط \lr{Gazebo} قابل دریافت بود. اما در سامانه واقعی، چنین مرجع مستقیمی وجود ندارد و موقعیت نسبی اسکوتر باید از داده‌های حسگرهای نصب‌شده روی سامانه تخمین زده شود. به همین دلیل، در پیاده‌سازی عملی، سرعت حرکتی از داده انکودرهای چرخ‌ها و جهت‌گیری از زاویه \lr{yaw} واحد اندازه‌گیری اینرسی استخراج شده و با ترکیب این داده‌ها، اودومتری نسبی اسکوتر نسبت به نقطه شروع آزمون محاسبه می‌شود. بنابراین، در حالی که در شبیه‌سازی موقعیت نسبی از وضعیت شبیه‌ساز به دست می‌آید، در سامانه واقعی این کمیت حاصل یک تخمین حسگری است و می‌تواند تحت تأثیر لغزش چرخ‌ها، نویز \lr{IMU} و خطای تجمعی انتگرال‌گیری قرار گیرد.

همچنین تفاوت میان مدل حرکتی شبیه‌سازی و سامانه فیزیکی واقعی باید در نظر گرفته شود. در شبیه‌سازی، پاسخ موتورها و اعمال فرمان به چرخ‌ها به‌صورت ایده‌آل‌تر و با محدودیت‌های کمتر مدل می‌شود، در حالی که در سامانه واقعی، ویژگی‌هایی مانند دینامیک موتور، تأخیر راه‌اندازها، محدودیت باتری، اشباع \lr{PWM}، اصطکاک، لقی مکانیکی و لغزش چرخ‌ها بر اجرای فرمان اثر می‌گذارند. به همین دلیل، در سامانه عملی فرمان خروجی عامل مستقیماً به موتورها اعمال نمی‌شود، بلکه ابتدا به سرعت هدف چهار موتور تبدیل شده و سپس آردوینو با استفاده از بازخورد انکودرها، کنترل پایین‌سطح سرعت موتورها را انجام می‌دهد. این لایه باعث می‌شود اثر بخشی از تفاوت‌های میان مدل شبیه‌سازی و سخت‌افزار واقعی کاهش یابد.

در بخش موانع نیز اصل کلی معماری حفظ شد. موانع ابتدا در تصویر واقعی تشخیص داده می‌شوند، فاصله آن‌ها با استفاده از نقشه اختلاف مکانی برآورد می‌شود و سپس موقعیت و سرعت نسبی آن‌ها با فیلتر کالمن پالایش می‌گردد. به این ترتیب، هم برای هدف و هم برای موانع، خروجی خام ادراک دیداری پیش از ورود به مشاهده عامل به داده‌ای پایدارتر و سازگارتر با ساختار آموزش تبدیل می‌شود.

جدول \ref{tab:sim_real_mapping_updated} نگاشت خلاصه اجزای اصلی معماری شبیه‌سازی به سامانه عملی را نشان می‌دهد.

\begin{table}[htb]
	\centering
	\caption{نگاشت خلاصه اجزای معماری شبیه‌سازی به سامانه عملی}
	\label{tab:sim_real_mapping_updated}
	\renewcommand{\arraystretch}{1.18}
	\begin{tabular}{|p{0.23\textwidth}|p{0.31\textwidth}|p{0.36\textwidth}|}
		\hline
		\textbf{بخش سامانه} &
		\textbf{در محیط شبیه‌سازی} &
		\textbf{در سامانه عملی} \\
		\hline
	
		ادراک دیداری و عمق &
		تصاویر استریو، تصحیح هندسی و تولید نقشه اختلاف مکانی با پارامترهای شبیه‌سازی &
		تصاویر دوربین واقعی، کالیبراسیون عملی، تصحیح هندسی و تولید نقشه اختلاف مکانی با همان منطق پردازشی \\
		\hline
		
		تشخیص و پالایش هدف &
		جداسازی هدف بر اساس ویژگی رنگی و استخراج فاصله، زاویه و سرعت نسبی هدف &
		تشخیص هدف با \lr{AprilTag}، استخراج فاصله و زاویه نسبی، و محاسبه سرعت نسبی هدف با فیلتر کالمن \\
		\hline
		
		تشخیص و پالایش موانع &
		تشخیص موانع، تخمین فاصله و پالایش زمانی داده‌ها در محیط شبیه‌سازی &
		تشخیص موانع از تصویر واقعی، تخمین عمق از نقشه اختلاف مکانی و پالایش موقعیت و سرعت نسبی موانع با فیلتر کالمن \\
		\hline
		
		وضعیت حرکتی و اودومتری &
		دریافت موقعیت، جهت‌گیری و سرعت از خروجی اودومتری و وضعیت \lr{Gazebo} &
		تخمین اودومتری نسبی با ترکیب سرعت حاصل از انکودرهای چرخ و زاویه \lr{yaw} واحد اندازه‌گیری اینرسی \\
		\hline
		
		اعمال فرمان به موتورها &
		اعمال فرمان به مدل حرکتی شبیه‌سازی‌شده با دینامیک ساده‌تر و محدودیت‌های کمتر &
		تبدیل فرمان عامل به سرعت هدف چهار موتور و اجرای آن با کنترل پایین‌سطح مبتنی بر بازخورد انکودر \\
		\hline
		
		حسگرهای نزدیک‌برد &
		مدل حسگرهای فراصوت برای مشاهده و سپر ایمنی &
		پنج حسگر فراصوت واقعی برای مشاهده عامل و تشخیص وضعیت‌های نزدیک‌برد \\
		\hline
		
		تشکیل مشاهده و اجرای مدل &
		ساخت مشاهده از داده‌های شبیه‌سازی و اجرای مدل آموزش‌دیده &
		تبدیل داده‌های واقعی به همان قالب مشاهده و اجرای همان مدل \lr{LSTM-SAC} در حالت استنتاج \\
		\hline
		
		فرمان و ایمنی &
		تولید فرمان حرکتی و اعمال منطق سپر ایمنی در محیط شبیه‌سازی &
		تبدیل فرمان عامل به فرمان قابل اجرا روی سخت‌افزار و اعمال محدودیت‌های ایمنی در مسیر اجرا \\
		\hline
		
	\end{tabular}

	
\end{table}

بر این اساس، سامانه عملی را می‌توان پیاده‌سازی فیزیکی معماری طراحی‌شده در شبیه‌سازی دانست. با این حال، اجرای واقعی صرفاً به انتقال مدل آموزش‌دیده محدود نیست؛ بلکه نیازمند زنجیره‌ای از واسط‌های ادراکی، ارتباطی و کنترلی است تا داده‌های خام حسگرها به مشاهده قابل استفاده برای عامل تبدیل شوند و فرمان خروجی عامل نیز به فرمان قابل اعمال روی موتورهای واقعی برسد. بنابراین، ارتباط میان شبیه‌سازی و سامانه عملی بر پایه حفظ منطق اصلی ادراک، مشاهده، تصمیم‌گیری، کنترل و ایمنی شکل گرفته است، در حالی که جزئیات اجرایی متناسب با محدودیت‌ها، نویزها و نرخ پردازش سخت‌افزار واقعی تنظیم شده‌اند.


\section{تجهیز و آماده‌سازی سخت‌افزاری سامانه عملی}
\label{sec:real_hardware_preparation}

پس از تعیین معماری کلی سامانه، نمونه فیزیکی اسکوتر برای نصب حسگرها، واحدهای پردازشی و اجزای کنترلی آماده شد. در این بخش، صرفاً نحوه جانمایی و نصب اجزای سخت‌افزاری اصلی توضیح داده می‌شود؛ جزئیات پردازش داده‌ها، تشکیل مشاهده و اجرای مدل تصمیم‌گیر در بخش‌های بعدی فصل ارائه خواهد شد.

\subsection{بستر مکانیکی اسکوتر و تغییرات اعمال‌شده}
\label{subsec:real_scooter_platform}

در این پژوهش، از یک اسکوتر چهارچرخ موجود به‌عنوان بستر مکانیکی استفاده شد. این اسکوتر از ابتدا دارای چهار موتور، راه‌اندازهای موتور و ساختار حرکتی مبتنی بر رانش لغزشی بود؛ بنابراین طراحی شاسی، انتخاب موتور و توسعه اولیه سامانه حرکتی در محدوده این پژوهش قرار نداشت. شاسی اصلی و مجموعه حرکتی اسکوتر بدون تغییر مکانیکی اساسی مورد استفاده قرار گرفتند و آماده‌سازی سخت‌افزاری عمدتاً به نصب حسگرها، جانمایی واحدهای پردازشی و کنترلی، و آماده‌سازی مسیرهای توان و سیگنال محدود شد.

فضای داخلی بدنه برای قرارگیری آردوینو مگا، راه‌اندازهای موتور، باتری‌ها و کابل‌های توان و سیگنال مورد استفاده قرار گرفت. در آزمون‌های آزمایشگاهی، لپ‌تاپ روی درپوش فوقانی اسکوتر قرار داده شد تا دسترسی به نرم‌افزار، ثبت داده‌ها و عیب‌یابی ساده‌تر باشد. این جانمایی برای نسخه آزمایشگاهی مناسب است، اما در یک نسخه نهایی می‌تواند با یک واحد پردازشی کوچک‌تر و قابل نصب دائم جایگزین شود.

\subsection{نصب دوربین استریو و رزبری پای}
\label{subsec:stereo_camera_raspberry_installation}

دوربین استریو در بخش جلویی اسکوتر نصب شد تا میدان دید آن با جهت حرکت اصلی اسکوتر هم‌راستا باشد. برای نصب دوربین، ابتدا یک پایه فلزی به ارتفاع ۷.۰ متر به قسمت جلویی بدنه پیچ شد. سپس یک صفحه از جنس \lr{ABS} به‌صورت موازی با پایه فلزی روی آن نصب گردید و ماژول دوربین استریو به این صفحه پیچ شد. این ساختار، سطح نصب مناسبی برای دوربین فراهم کرد و موقعیت آن را نسبت به بدنه ثابت نگه داشت.

رزبری پای ۵ نیز در نزدیکی دوربین و روی همان مجموعه نصب جلویی قرار گرفت. برای این منظور، در پشت پایه فلزی، یک صفحه \lr{ABS} دیگر به‌صورت عمودی و با استفاده از براکت به پایه متصل شد و رزبری پای روی این صفحه پیچ گردید. فاصله رزبری پای تا دوربین در حدود ۵ سانتی‌متر است؛ بنابراین اتصال دو دوربین به رزبری از طریق کابل‌های \lr{CSI} بدون کابل‌کشی بلند انجام می‌شود.
%عکس از پایه و سازه متصل شده%

\subsection{نصب حسگرهای فراصوت}
\label{subsec:ultrasonic_sensor_installation}

پنج حسگر فراصوت \lr{SRF05} در موقعیت‌های جلو-چپ، جلو-راست، چپ، راست و عقب در ارتفاع حدودی از زمین ۵۱.۰ متر نسب به زمین نصب شدند. محل نصب این حسگرها مطابق نواحی تعریف‌شده در فصل چهارم انتخاب شد تا داده‌های واقعی آن‌ها با منطق فضای مشاهده و سپر ایمنی سازگار باشد.

برای نصب هر حسگر، یک نگهدارنده اختصاصی طراحی شد. این نگهدارنده دارای دو خار بود که در حفره‌های موجود روی برد حسگر فراصوت قرار می‌گرفتند و حسگر را نسبت به نگهدارنده ثابت می‌کردند. سپس مجموعه حسگر و نگهدارنده در نقاط تعیین‌شده روی بدنه نصب شد. در هنگام نصب، جهت‌گیری هر حسگر به‌گونه‌ای تنظیم شد که ناحیه مورد نظر خود را پوشش دهد و توسط بدنه یا تجهیزات جانبی مسدود نشود.
%عکس از قاب و سنسور وصل شده%

\subsection{جانمایی واحد اندازه‌گیری اینرسی، آردوینو، راه‌اندازهای موتور و سامانه تغذیه}
\label{subsec:imu_arduino_driver_power_installation}

واحد اندازه‌گیری اینرسی از طریق اتصال \lr{USB} به رزبری پای متصل شد. آردوینو در داخل محفظه اصلی بدنه نصب شد و ارتباط با انکودرهای چهار موتور، حسگرهای فراصوت و راه‌اندازهای موتور را بر عهده گرفت. راه‌اندازهای موتور نیز داخل بدنه و روی یک صفحه فلزی ثابت شدند. این صفحه علاوه بر ایجاد محل نصب مکانیکی، به دفع حرارت تولیدشده توسط راه‌اندازها کمک می‌کند. راه‌اندازها سیگنال کنترلی را از آردوینو دریافت کرده و توان لازم را به موتورهای چرخ منتقل می‌کنند.

سامانه تغذیه از چهار باتری لیتیوم‌یونی \lr{Pleson 10S2P} تشکیل شده است. مشخصات هر باتری شامل ولتاژ نامی $36~V$، ظرفیت $4.4~Ah$، انرژی $158.4~Wh$ و ولتاژ شارژ $42~V$ است. باتری‌ها به‌صورت موازی بسته شده‌اند و باس توان $36~V$ را تشکیل می‌دهند. این باس برای تغذیه راه‌اندازهای موتور استفاده می‌شود و از طریق تنظیم ولتاژ، تغذیه $5~V$ موردنیاز رزبری پای، آردوینو و حسگرهای وابسته نیز فراهم می‌گردد. لپ‌تاپ در نسخه آزمایشگاهی از منبع تغذیه مستقل خود استفاده می‌کند.

شکل \ref{fig:power_and_signal_distribution} شماتیک مسیرهای توان و سیگنال در سامانه عملی را نشان می‌دهد.

\begin{figure}[htbp]
	\centering
	\includegraphics[width=0.75\linewidth]{schematic_of_electronic_power_system}
	\caption{شماتیک مسیرهای توان و سیگنال در سامانه عملی اسکوتر خودران}
	\label{fig:power_and_signal_distribution}
\end{figure}

پس از تکمیل جانمایی اجزای مکانیکی، حسگرها، واحدهای پردازشی، باتری‌ها
و راه‌اندازهای موتور، مدل سه‌بعدی یکپارچه سامانه برای بررسی وضعیت
نهایی مونتاژ تهیه شد. شکل \ref{fig:assembled_scooter_overview}
نمای کلی اسکوتر تجهیزشده را از چند زاویه مختلف نشان می‌دهد. این
نماها امکان بررسی هم‌زمان ساختار شاسی، مجموعه تعلیق و چرخ‌ها،
جانمایی باتری‌ها و تجهیزات الکترونیکی داخل بدنه، و نحوه نصب مجموعه
دوربین استریو و رزبری پای روی پایه جلویی را فراهم می‌کنند. نمای
بالا نیز توزیع اجزای داخلی و استفاده از فضای موجود در محفظه اصلی
اسکوتر را با وضوح بیشتری نشان می‌دهد.

\begin{figure}[htb]
	\centering
	
	\subfloat[]{
		\includegraphics[width=0.47\textwidth]
		{Scooter_3d_3corner_view_1.png}
	}
	\hfill
	\subfloat[]{
		\includegraphics[width=0.47\textwidth]
		{Scooter_3d_3corners_view_2.png}
	}
	
	\vspace{0.3cm}
	
	\subfloat[]{
		\includegraphics[width=0.47\textwidth]
		{Scooter_3d_rear_view.png}
	}
	\hfill
	\subfloat[]{
		\includegraphics[width=0.47\textwidth]
		{Scooter_3d_upper_view.png}
	}
	
	\caption{نمای کلی مدل سه‌بعدی مونتاژشده اسکوتر و جانمایی اجزای
		سخت‌افزاری: الف) نمای سه‌بعدی از سمت جلو و راست،
		ب) نمای سه‌بعدی از سمت جلو و چپ،
		ج) نمای روبه‌رو از مجموعه تعلیق، چرخ‌ها و پایه تجهیزات،
		د) نمای بالا از جانمایی باتری‌ها، راه‌اندازهای موتور،
		واحدهای کنترلی و تجهیزات پردازشی}
	\label{fig:assembled_scooter_overview}
\end{figure}

\section{پیاده‌سازی لایه ادراک در سامانه عملی}
\label{sec:real_perception_layer}

پس از آماده‌سازی سخت‌افزاری سامانه، داده‌های خام حسگرها باید به اطلاعات قابل استفاده برای ناوبری تبدیل شوند. لایه ادراک در سامانه عملی وظیفه دریافت و پردازش تصاویر استریو، داده‌های حسگرهای فراصوت، اطلاعات انکودرها و داده واحد اندازه‌گیری اینرسی را بر عهده دارد. خروجی این لایه شامل اطلاعاتی مانند فاصله و زاویه هدف، موقعیت نسبی موانع، فاصله‌های نزدیک‌برد و وضعیت نسبی حرکت اسکوتر است که در ادامه برای تشکیل مشاهده عامل تصمیم‌گیر استفاده می‌شوند.

در پیاده‌سازی عملی، تلاش شد منطق کلی لایه ادراک با ساختار تعریف‌شده در فصل چهارم سازگار باقی بماند. با این حال، به دلیل استفاده از حسگرهای واقعی، پارامترهای کالیبراسیون، نرخ پردازش، شرایط نوری و نویز اندازه‌گیری با محیط شبیه‌سازی متفاوت هستند. همچنین در تشخیص هدف، به‌جای جداسازی رنگی مورد استفاده در شبیه‌سازی، از نشانگر \lr{AprilTag} استفاده شد؛ زیرا تصاویر استریو در سامانه عملی به‌صورت خاکستری منتقل می‌شوند و اطلاعات رنگی برای تشخیص هدف در دسترس نیست.

\subsection{دریافت تصویر از دوربین استریو}
\label{subsec:real_stereo_image_acquisition}

برای دریافت اطلاعات دیداری محیط، از یک دوربین استریو مبتنی بر دو حسگر \lr{IMX219-83} استفاده شد. این دو حسگر به‌صورت چپ و راست عمل کرده و از طریق دو کابل \lr{CSI} به رزبری متصل هستند. در سامانه عملی، هر دو تصویر ابتدا روی رزبری دریافت می‌شوند و سپس برای پردازش‌های بعدی به لپ‌تاپ منتقل می‌گردند. هدف از این مرحله، فراهم‌سازی یک جفت تصویر هم‌زمان از محیط روبه‌روی اسکوتر است تا در ادامه بتوان از آن‌ها برای تصحیح هندسی، محاسبه اختلاف مکانی، تخمین عمق، تشخیص هدف و تشخیص موانع استفاده کرد.

تصاویر دریافتی از دوربین با وضوح ($1640 \times 1232$) ثبت می‌شوند و نرخ انتقال آن‌ها در پیاده‌سازی عملی حدود ۱۵ فریم بر ثانیه است. انتقال تصاویر از رزبری پای به لپ‌تاپ از طریق کابل \lr{LAN} و با استفاده از \lr{GStreamer} انجام می‌شود. انتخاب این ساختار انتقالی باعث شد دریافت تصاویر چپ و راست با پایداری مناسب انجام شود و در عین حال، بار ارتباطی و پردازشی سامانه در سطح قابل قبولی باقی بماند.

در پیاده‌سازی عملی، فریم‌های خام دوربین در قالب \lr{YUV420} دریافت می‌شوند. در این قالب، هر تصویر از سه مؤلفه اصلی تشکیل شده است: صفحه \lr{Y} که شدت روشنایی را نگه می‌دارد، و دو صفحه \lr{U} و \lr{V} که اطلاعات رنگی را بیان می‌کنند. همان‌گونه که در شکل \ref{fig:real_stereo_image_pipeline} نشان داده شده است، در این پژوهش تنها صفحه \lr{Y} از هر تصویر استخراج و برای ادامه پردازش نگه‌داری می‌شود و صفحات \lr{U} و \lr{V} حذف می‌شوند. دلیل این انتخاب آن است که در پردازش استریو، اطلاعات روشنایی برای تطبیق نقاط متناظر و برآورد اختلاف مکانی کافی است و استفاده از صفحات رنگی نه‌تنها ضرورت مستقیمی ندارد، بلکه باعث افزایش حجم داده و پیچیدگی پردازش نیز می‌شود. در نتیجه، خروجی این مرحله برای هر دوربین یک تصویر خاکستری ۸ بیتی است که ورودی اصلی زنجیره پردازش استریو را تشکیل می‌دهد.

استفاده از صفحه \lr{Y} دو مزیت اصلی دارد. نخست، حجم داده ارسالی و پردازش‌شده نسبت به استفاده از تصویر رنگی کامل کاهش می‌یابد و این موضوع برای اجرای سامانه در زمان واقعی اهمیت دارد. دوم، مرحله تطبیق استریو بر پایه الگوهای شدت روشنایی انجام می‌شود و بنابراین استفاده از تصویر خاکستری برای تولید نقشه اختلاف مکانی کفایت می‌کند. به همین دلیل، در پیاده‌سازی عملی، هر دو تصویر چپ و راست پس از استخراج صفحه \lr{Y}، به‌عنوان ورودی مشترک مراحل بعدی شامل تصحیح هندسی، تطبیق استریو و برآورد عمق مورد استفاده قرار می‌گیرند.

بر این اساس، مرحله دریافت تصویر در سامانه عملی تنها یک فرآیند انتقال ساده داده نیست، بلکه شامل دریافت جفت‌تصویر استریو، تبدیل فریم خام \lr{YUV420} به داده روشنایی قابل استفاده، و آماده‌سازی ورودی مناسب برای پردازش‌های بعدی است. این مرحله، نقطه شروع کل زنجیره ادراک دیداری در سامانه عملی محسوب می‌شود.

\begin{figure}[htbp]
	\centering
	\includegraphics[width=0.75\linewidth]{stereo_vision_y_plane_extraction_process.png}
	\caption{روند دریافت و آماده‌سازی تصاویر دوربین استریو در سامانه عملی؛ از دریافت فریم خام تا استخراج صفحه روشنایی \lr{Y} و استفاده از آن در پردازش استریو}
	\label{fig:real_stereo_image_pipeline}
\end{figure}

\subsection{تصحیح هندسی تصاویر و تولید نقشه اختلاف مکانی}
\label{subsec:real_depth_map_generation}

پس از دریافت تصاویر چپ و راست در لپ‌تاپ، ابتدا تصحیح هندسی تصاویر استریو انجام می‌شود. هدف از این مرحله، حذف اثر اعوجاج لنزها و هم‌راستاسازی تصاویر چپ و راست است؛ به‌گونه‌ای که نقاط متناظر در دو تصویر روی خطوط افقی متناظر قرار گیرند. این هم‌راستاسازی برای محاسبه پایدار اختلاف مکانی ضروری است، زیرا در صورت وجود خطای هندسی میان دو تصویر، الگوریتم تطبیق استریو نمی‌تواند جست‌وجوی نقاط متناظر را به‌صورت یک‌بعدی و پایدار انجام دهد.

تصحیح هندسی با استفاده از پارامترهای کالیبراسیون دوربین استریو انجام شد. در کالیبراسیون نهایی، رزولوشن تصویر برابر ($1640 \times 1232$) پیکسل، فاصله پایه دوربین‌ها برابر ($60.172~mm$)، فاصله کانونی تصویر تصحیح‌شده برابر ($1238.075~px$)، و خطای بازفرافکنی استریو برابر ($0.3774~px$) به‌دست آمد. این مقادیر برای ساخت نگاشت‌های تصحیح و همچنین محاسبات فاصله بر اساس اختلاف مکانی مورد استفاده قرار گرفتند. جزئیات کامل روند کالیبراسیون، ماتریس‌های درونی دوربین‌ها، ضرایب اعوجاج، ماتریس‌های تصحیح و ماتریس‌های تصویر در پیوست \ref{app:stereo_calibration_details} و جدول \ref{tab:stereo_calibration_final_params} ارائه شده‌اند.

در ساختار نرم‌افزاری سامانه، تصحیح هندسی فقط در گره دریافت تصویر انجام می‌شود. این گره پس از دریافت تصاویر چپ و راست، فایل کالیبراسیون را بارگذاری کرده، نگاشت‌های تصحیح را می‌سازد و تصاویر تصحیح‌شده را روی موضوعات ارتباطی مربوط به تصویر چپ و راست منتشر می‌کند. پس از انتشار تصاویر تصحیح‌شده، گره تولید اختلاف مکانی، تصاویر چپ و راست را دریافت کرده و با استفاده از هم‌زمان‌سازی تقریبی، جفت‌تصویرهای نزدیک از نظر زمانی را برای پردازش انتخاب می‌کند. سپس تصاویر خاکستری تصحیح‌شده وارد الگوریتم \lr{StereoSGBM} می‌شوند. این الگوریتم برای هر ناحیه از تصویر چپ، ناحیه متناظر آن را در تصویر راست جست‌وجو کرده و اختلاف افقی میان آن‌ها را به‌عنوان مقدار اختلاف مکانی محاسبه می‌کند. مبنای هندسی تبدیل اختلاف مکانی به عمق همان رابطه ارائه‌شده در فصل چهارم است و در این بخش از تکرار آن خودداری می‌شود؛ برای رابطه عمق بر حسب اختلاف مکانی به رابطه \ref{eq:stereo_depth_basic} مراجعه شود.

در پیاده‌سازی عملی، پارامترهای \lr{SGBM} بر اساس رزولوشن تصویر و محدوده فاصله مورد نیاز در آزمون‌ها تنظیم شدند. مقدار حداقل اختلاف مکانی برابر ۱۶ و مقدار بازه اختلاف مکانی برابر ۱۷۶ در نظر گرفته شد تا بازه مناسبی از اختلاف مکانی برای فاصله‌های کاری سامانه پوشش داده شود. همچنین اندازه بلوک برابر ۷ انتخاب شد تا میان کاهش نویز و حفظ مرزهای عمقی تعادل برقرار شود. از آنجا که محاسبه اختلاف مکانی روی تصویر با وضوح بالا هزینه پردازشی زیادی دارد، در پیاده‌سازی عملی امکان کاهش مقیاس داخلی تصویر پیش از اجرای \lr{SGBM} در نظر گرفته شد. پس از محاسبه اختلاف مکانی در مقیاس کوچک‌تر، مقادیر اختلاف مکانی به واحد پیکسل در وضوح اصلی بازگردانده می‌شوند.

برای بهبود کیفیت خروجی، دو مرحله تکمیلی نیز در زنجیره پردازش قرار گرفت. نخست، از روش \lr{CLAHE} برای بهبود کنتراست موضعی تصاویر خاکستری استفاده شد. این پیش‌پردازش در نواحی کم‌بافت یا دارای تغییرات روشنایی ضعیف می‌تواند به افزایش قابلیت تطبیق میان دو تصویر کمک کند. دوم، از فیلتر \lr{WLS} برای پالایش نقشه اختلاف مکانی استفاده شد. نقشه اختلاف مکانی خام معمولاً شامل نویز، لکه‌های پراکنده و ناپایداری در مرز اجسام است؛ فیلتر \lr{WLS} با استفاده از تصویر مرجع، خروجی تطبیق را هموارتر کرده و در عین حال تلاش می‌کند مرزهای مهم تصویر حفظ شوند. جزئیات کامل الگوریتم تولید نقشه اختلاف مکانی، تنظیمات \lr{SGBM}، نقش \lr{WLS}، اثر \lr{CLAHE} و نحوه تبدیل اختلاف مکانی به عمق در پیوست \ref{appendix:stereo_depth_sgbm} ارائه شده است.

خروجی این مرحله، نقشه اختلاف مکانی متراکم در قالب پیام \lr{DisparityImage} است. در این پیام، علاوه بر تصویر اختلاف مکانی، پارامترهای مورد نیاز برای تفسیر مقادیر اختلاف مکانی، از جمله فاصله کانونی و فاصله پایه دوربین، نیز قرار می‌گیرند. این نقشه به‌صورت مستقیم وارد فضای مشاهده عامل نمی‌شود، بلکه در مراحل بعدی برای استخراج فاصله هدف و موانع مورد استفاده قرار می‌گیرد.

شکل \ref{fig:real_disparity_generation_process} نمونه‌ای از مراحل اصلی تولید و پالایش نقشه اختلاف مکانی در سامانه عملی را نشان می‌دهد. در این شکل، ابتدا تصویر ورودی خام نمایش داده شده است. سپس نقشه اختلاف مکانی خام حاصل از اجرای \lr{SGBM}، نقشه پالایش‌شده با \lr{WLS} و خروجی حاصل از پیش‌پردازش \lr{CLAHE} ارائه شده‌اند. این مقایسه نشان می‌دهد که کیفیت نقشه اختلاف مکانی در سامانه واقعی به کیفیت تصویر ورودی، کنتراست موضعی و پالایش پس از تطبیق وابسته است.

\begin{figure}[htbp]
	\centering
	
	\subfloat[]{
		\includegraphics[width=0.48\linewidth]{real_stereo_raw_input.png}
	}
	\hfill
	\subfloat[]{
		\includegraphics[width=0.48\linewidth]{real_disparity_raw.png}
	}
	
	\vspace{0.4cm}
	
	\subfloat[]{
		\includegraphics[width=0.48\linewidth]{real_disparity_wls.png}
	}
	\hfill
	\subfloat[]{
		\includegraphics[width=0.48\linewidth]{real_disparity_clahe.png}
	}
	
	\caption{نمونه مراحل تولید و پالایش نقشه اختلاف مکانی در سامانه عملی:
		الف) تصویر ورودی خام،
		ب) نقشه اختلاف مکانی خام،
		ج) نقشه اختلاف مکانی پس از اعمال \lr{WLS}،
		د) نقشه اختلاف مکانی با پیش‌پردازش \lr{CLAHE}}
	\label{fig:real_disparity_generation_process}
\end{figure}

\subsection{تشخیص هدف با استفاده از نشانگر \lr{AprilTag}}
\label{subsec:real_apriltag_goal_detection}

در محیط شبیه‌سازی، هدف بر اساس ویژگی رنگی آن از تصویر جدا می‌شد. اما در سامانه عملی، همان‌طور که در بخش‌های پیشین توضیح داده شد، تصاویر استریو به‌صورت خاکستری منتقل می‌شوند و اطلاعات رنگی برای تشخیص هدف در دسترس نیست. به همین دلیل، در پیاده‌سازی واقعی از نشانگرهای \lr{AprilTag} برای تشخیص هدف استفاده شد. این روش نسبت به جداسازی رنگی وابستگی کمتری به رنگ هدف دارد و در تصویر خاکستری نیز قابل استفاده است.

برای هدف آزمایشگاهی، پنج نشانگر \lr{AprilTag} که در شکل \ref{fig:April_Tags} قابل مشاهده هستند، روی سطوح مختلف جعبه هدف نصب شد. این نشانگرها روی سطوح جانبی و سطح بالایی جعبه قرار گرفتند و روی سطح زیرین جعبه نشانگری نصب نشد. دو سطح حامل تگ دارای ابعاد تقریبی ($50 \times 60~cm$)، دو سطح دیگر دارای ابعاد تقریبی ($30 \times 50~cm$)، و سطح بالایی دارای ابعاد تقریبی ($30 \times 60~cm$) بود. استفاده از چند تگ روی وجوه مختلف هدف باعث شد در زوایا و وضعیت‌های مختلف، حداقل یکی از نشانگرها در میدان دید دوربین قرار گیرد و امکان تشخیص هدف افزایش یابد.

\begin{figure}[htbp]
	\centering
	\includegraphics[width=0.75\linewidth]{April_Tag_ID_0-->4}
	\caption{نشانگرهای استفاده‌شده روی هدف}
	\label{fig:April_Tags}
\end{figure}

خانواده نشانگرهای مورد استفاده از نوع \lr{tag36h11} است و شناسه‌های مجاز برای تشخیص هدف در بازه تگ‌های نصب‌شده تعریف شده‌اند. در صورتی که چند نشانگر به‌صورت هم‌زمان در تصویر دیده شوند، نشانگری که سطح تصویر بزرگ‌تری دارد به‌عنوان نماینده هدف انتخاب می‌شود. این انتخاب از آن جهت منطقی است که معمولاً تگ بزرگ‌تر در تصویر، واضح‌تر و پایدارتر قابل تشخیص است و مقدار اختلاف مکانی استخراج‌شده از ناحیه آن نیز قابلیت اطمینان بیشتری دارد.

پس از تشخیص تگ، چهار گوشه آن در تصویر استخراج شده و مرکز تگ محاسبه می‌شود. مرکز تگ برای تعیین زاویه نسبی هدف نسبت به محور دید دوربین به کار می‌رود. به بیان دیگر، اختلاف افقی مرکز تگ نسبت به مرکز تصویر، همراه با فاصله کانونی دوربین، برای محاسبه زاویه دید هدف استفاده می‌شود. خروجی اولیه این مرحله شامل مختصات مرکز تگ در تصویر، زاویه نسبی هدف و شناسه تگ تشخیص‌داده‌شده است.

برای محاسبه فاصله هدف، پس از تعیین کادر محدودکننده تگ در تصویر چپ، ناحیه متناظر آن روی نقشه اختلاف مکانی انتخاب می‌شود. سپس مقادیر معتبر اختلاف مکانی در این ناحیه استخراج شده و با استفاده از پارامترهای فاصله کانونی و فاصله پایه دوربین، عمق هدف محاسبه می‌شود. مبنای محاسبه عمق همان رابطه اختلاف مکانی ارائه‌شده در فصل چهارم و بخش \ref{subsec:real_depth_map_generation} است.

برای افزایش پایداری تخمین فاصله، کادر تگ به‌صورت محدود گسترش داده می‌شود تا تعداد پیکسل‌های معتبر بیشتری از نقشه اختلاف مکانی در محاسبه استفاده شوند. سپس مقادیر نامعتبر یا کمتر از حد مجاز حذف می‌شوند و در صورتی که تعداد پیکسل‌های معتبر کافی باشد، مقدار نماینده اختلاف مکانی از روی آن ناحیه محاسبه می‌شود. در پیاده‌سازی انجام‌شده، مقدار میانه به‌عنوان معیار اصلی انتخاب شد، زیرا نسبت به نقاط پرت ناشی از خطای تطبیق استریو، مرزهای تگ و نواحی نامعتبر مقاوم‌تر از میانگین ساده است.

پس از استخراج فاصله و زاویه نسبی هدف، مختصات نسبی هدف در دستگاه مختصات اسکوتر محاسبه می‌شود. سپس برای کاهش نوسان اندازه‌گیری و تخمین سرعت نسبی هدف، از همان ساختار فیلتر کالمن استفاده می‌شود که برای پالایش موقعیت و سرعت موانع به کار رفته است. در این روش، بردار حالت شامل موقعیت طولی، سرعت طولی، موقعیت عرضی و سرعت عرضی هدف نسبت به اسکوتر است و اندازه‌گیری ورودی فیلتر از مختصات دوبعدی هدف تشکیل می‌شود. بنابراین، فیلتر کالمن علاوه بر هموارسازی موقعیت نسبی هدف، مؤلفه‌های سرعت نسبی آن را نیز تخمین می‌زند.

استفاده از فیلتر کالمن برای هدف باعث می‌شود تغییرات لحظه‌ای ناشی از نویز نقشه اختلاف مکانی، جابه‌جایی مرکز تگ در تصویر و ناپایداری تشخیص در چند فریم متوالی کاهش یابد. این موضوع به‌ویژه برای تشکیل مشاهده عامل اهمیت دارد، زیرا عامل تصمیم‌گیر علاوه بر فاصله و زاویه هدف، به مؤلفه‌های سرعت نسبی هدف نیز نیاز دارد. جزئیات مدل حالت، مدل اندازه‌گیری و مراحل پیش‌بینی و تصحیح فیلتر کالمن در پیوست \ref{appendix:kalman_speed} ارائه شده است.

در صورت عدم مشاهده هیچ‌یک از تگ‌های مجاز، خروجی گره تشخیص هدف با پرچم عدم مشاهده منتشر می‌شود. در این حالت، مقدار قابل اعتماد جدیدی برای فاصله، زاویه و سرعت نسبی هدف تولید نمی‌شود. در سطح گره تشخیص، وضعیت هدف به‌صورت نامرئی گزارش می‌شود و مختصات تصویر نیز مقدار نامعتبر می‌گیرد. سپس در مرحله تشکیل مشاهده عامل، از این پرچم برای مشخص کردن وضعیت عدم مشاهده هدف استفاده می‌شود و در صورت نیاز، آخرین موقعیت معتبر مشاهده‌شده برای حفظ پیوستگی محاسبات مورد استفاده قرار می‌گیرد. جزئیات این رفتار در بخش مربوط به تشکیل مشاهده عامل توضیح داده خواهد شد.

خروجی نهایی این گره در قالب وضعیت هدف منتشر می‌شود و شامل قابل مشاهده بودن هدف، فاصله تخمین‌زده‌شده، زاویه نسبی هدف، مختصات مرکز تگ در تصویر و مؤلفه‌های سرعت نسبی هدف است. این خروجی یکی از ورودی‌های اصلی زنجیره تشکیل مشاهده عامل محسوب می‌شود، زیرا عامل تصمیم‌گیر برای ناوبری هدف‌محور به فاصله، زاویه و سرعت نسبی هدف نیاز دارد.

%\begin{figure}[htbp]
%	\centering
%	\includegraphics[width=0.85\linewidth]{apriltag_goal_detection_result.png}
%	\caption{نمونه تشخیص هدف با استفاده از نشانگر \lr{AprilTag} روی تصویر چپ تصحیح‌شده}
%	\label{fig:apriltag_goal_detection_result}
%\end{figure}

\subsection{تشخیص موانع و استخراج موقعیت نسبی}
\label{subsec:real_obstacle_detection}

پس از تولید نقشه اختلاف مکانی، مرحله بعدی در لایه ادراک، تشخیص موانع و استخراج موقعیت نسبی آن‌ها نسبت به اسکوتر است. در سامانه عملی، این کار با ترکیب خروجی آشکارساز شیء و نقشه اختلاف مکانی انجام شد. به این ترتیب، آشکارساز، محل تقریبی مانع را در تصویر مشخص می‌کند و نقشه اختلاف مکانی، اطلاعات فاصله‌ای همان ناحیه را فراهم می‌سازد.

روند کلی این پردازش در شکل \ref{fig:real_obstacle_detection_depth_pipeline} نشان داده شده است. مطابق این شکل، تصاویر چپ و راست تصحیح‌شده از یک سو وارد گره تولید نقشه اختلاف مکانی می‌شوند و از سوی دیگر، تصویر چپ تصحیح‌شده برای تشخیص موانع به گره \lr{YOLO} داده می‌شود. سپس خروجی این دو مسیر، یعنی نقشه اختلاف مکانی و کادرهای محدودکننده موانع، در گره تخمین عمق ترکیب می‌شوند تا فاصله، زاویه و موقعیت نسبی هر مانع محاسبه شود.

در پیاده‌سازی انجام‌شده، تشخیص موانع فقط روی تصویر چپ تصحیح‌شده انجام می‌شود. از آنجا که تصاویر ورودی سامانه به‌صورت تک‌رنگ هستند، تصویر خاکستری ابتدا به قالب سه‌کاناله تبدیل می‌شود تا با ورودی مورد انتظار مدل \lr{YOLO} سازگار باشد. سپس تصویر برای کاهش هزینه پردازشی به اندازه ورودی مدل تغییر مقیاس داده می‌شود. پس از اجرای آشکارسازی، مختصات کادرهای محدودکننده دوباره به اندازه اصلی تصویر بازگردانده می‌شوند تا با نقشه اختلاف مکانی هم‌خوانی داشته باشند.

خروجی گره \lr{YOLO} شامل کادر محدودکننده، کلاس مانع، میزان اطمینان تشخیص و شناسه مانع است. در این پژوهش، از این خروجی برای تعیین ناحیه متناظر هر مانع روی نقشه اختلاف مکانی استفاده شد. به بیان دیگر، \lr{YOLO} وظیفه تشخیص محل مانع در تصویر را بر عهده دارد، اما فاصله مانع مستقیماً از مدل تشخیص به دست نمی‌آید؛ بلکه با نمونه‌برداری از نقشه اختلاف مکانی در محدوده همان کادر محاسبه می‌شود.

برای محاسبه عمق هر مانع، ابتدا کادر محدودکننده آن روی نقشه اختلاف مکانی قرار داده می‌شود. سپس ناحیه داخلی کادر تا حدی کوچک‌تر می‌شود تا اثر پیکسل‌های مرزی، نواحی انسداد و خطاهای تطبیق در اطراف جسم کاهش یابد. پس از آن، پیکسل‌های نامعتبر یا دارای اختلاف مکانی نامناسب حذف می‌شوند. اگر تعداد پیکسل‌های معتبر از حداقل تعیین‌شده بیشتر باشد، یک مقدار نماینده برای اختلاف مکانی مانع محاسبه می‌شود. در حالت پیش‌فرض، از میانه مقادیر معتبر برای گزارش مقدار اختلاف مکانی استفاده شده است.
پس از محاسبه اختلاف مکانی نماینده، فاصله مانع با استفاده از پارامترهای کالیبراسیون دوربین و رابطه عمق-اختلاف مکانی که در فصل چهارم معرفی شد، به دست می‌آید. زاویه نسبی مانع نیز از روی موقعیت افقی مرکز کادر نسبت به مرکز تصویر محاسبه می‌شود. با داشتن فاصله و زاویه، مختصات نسبی مانع در دستگاه مختصات اسکوتر به دست می‌آید. در این دستگاه، مؤلفه طولی نشان‌دهنده فاصله مانع در راستای روبه‌روی اسکوتر و مؤلفه عرضی نشان‌دهنده انحراف جانبی مانع نسبت به محور حرکت است.

خروجی این مرحله برای هر مانع شامل شناسه، کلاس، فاصله پالایش‌شده، فاصله خام، مختصات نسبی، زاویه نسبی و اطلاعات سرعت تخمینی است. این اطلاعات در قالب یک پیام متنی ساختاریافته منتشر می‌شوند و در مراحل بعدی برای تشکیل مشاهده عامل و اعمال منطق ایمنی مورد استفاده قرار می‌گیرند. تخمین سرعت نسبی و پالایش زمانی موقعیت موانع در مرحله بعد و با استفاده از فیلتر کالمن انجام می‌شود که جزئیات آن در بخش \ref{subsec:real_obstacle_kalman_filtering} و پیوست \ref{appendix:kalman_speed} ارائه شده است.

\subsection{پالایش و ردیابی موانع با فیلتر کالمن}
\label{subsec:real_obstacle_kalman_filtering}

خروجی خام حاصل از ترکیب آشکارسازی \lr{YOLO} و نقشه اختلاف مکانی، به دلیل نویز تصویر، خطای تطبیق استریو، تغییرات کادر تشخیص و ناپایداری پیکسل‌های معتبر، دارای نوسان است. برای کاهش این نوسان و فراهم‌سازی داده پایدارتر برای عامل تصمیم‌گیر، موقعیت نسبی موانع با استفاده از فیلتر کالمن پالایش شد. مبانی ریاضی فیلتر کالمن، مدل حرکت سرعت ثابت، مدل اندازه‌گیری، پالایش مقادیر پرت عمق و نحوه تخمین سرعت نسبی موانع در پیوست \ref{appendix:kalman_speed} ارائه شده است؛ بنابراین در این بخش تنها نحوه به‌کارگیری آن در سامانه عملی توضیح داده می‌شود.

در پیاده‌سازی عملی، برای هر مانع معتبر یک فیلتر کالمن مستقل تعریف شد. بردار حالت هر مانع شامل موقعیت طولی، سرعت طولی، موقعیت عرضی و سرعت عرضی نسبت به دستگاه مختصات اسکوتر است. اندازه‌گیری ورودی فیلتر نیز از موقعیت دوبعدی مانع تشکیل می‌شود که از ترکیب عمق تخمین‌زده‌شده و زاویه نسبی مانع به دست می‌آید. بنابراین، فیلتر کالمن در این بخش فقط برای هموارسازی عمق استفاده نشده، بلکه برای پالایش موقعیت نسبی دوبعدی و تخمین سرعت نسبی مانع به کار رفته است.

پیش از به‌روزرسانی فیلتر کالمن، عمق خام هر مانع در یک بافر کوتاه‌مدت ذخیره شده و مقادیر پرت آن با استفاده از پالایش مبتنی بر بازه بین‌چارکی \lr{IQR} کاهش داده می‌شود. این پالایش روی مقدار عمق انجام می‌شود، نه روی اختلاف مکانی، زیرا رابطه میان عمق و اختلاف مکانی غیرخطی است. سپس مختصات نسبی مانع با استفاده از عمق پالایش‌شده و زاویه نسبی محاسبه شده و به‌عنوان اندازه‌گیری ورودی فیلتر مورد استفاده قرار می‌گیرد.

در هر گام زمانی، ابتدا بازه زمانی میان دو مشاهده متوالی محاسبه شده و در ماتریس گذار حالت اعمال می‌شود. سپس مرحله پیش‌بینی و اصلاح فیلتر کالمن انجام می‌گیرد. خروجی نهایی شامل موقعیت پالایش‌شده مانع، عمق پالایش‌شده، زاویه نسبی، مؤلفه‌های سرعت نسبی و همچنین مقدار خام عمق برای مقایسه و عیب‌یابی است. باید توجه داشت که سرعت‌های تخمین‌زده‌شده در این مرحله سرعت نسبی نسبت به اسکوتر هستند و به‌عنوان سرعت مطلق مانع در دستگاه مختصات جهانی تفسیر نمی‌شوند.

پارامترهای فیلتر در پیاده‌سازی عملی با هدف ایجاد تعادل میان نرمی خروجی و پاسخ‌گویی به حرکت موانع تنظیم شدند. در این تنظیمات، نویز فرایند موقعیت و سرعت به‌ترتیب برابر ($0.10$) و ($0.80$)، نویز اندازه‌گیری برابر ($0.10$)، و طول بافر عمق برابر ۸ نمونه در نظر گرفته شد. این مقادیر باعث شدند خروجی نسبت به جهش‌های لحظه‌ای عمق خام مقاوم‌تر شود، بدون آنکه تأخیر زیادی در دنبال‌کردن تغییرات موقعیت مانع ایجاد گردد.

\begin{figure}[htb]
	\centering
	\includegraphics[width=0.95\linewidth]{stereo_vision_pipeline_diagram}
	\caption{زنجیره پردازش تصویر استریو، تشخیص موانع و تخمین عمق و موقعیت نسبی موانع در سامانه عملی}
	\label{fig:real_obstacle_detection_depth_pipeline}
\end{figure}



\subsection{جمع‌بندی بخش بینایی لایه ادراک سامانه عملی}
\label{subsec:real_perception_layer_summary}

\begin{figure}[htbp]
	\centering
	\includegraphics[width=0.75\linewidth]{stereo_image_acquisition_pipeline_diagram.png}
	\caption{جمع‌بندی زنجیره لایه ادراک در سامانه عملی، از دریافت تصویر استریو تا تولید خروجی‌های مورد نیاز برای برآورد عمق، تشخیص موانع و تشخیص هدف}
	\label{fig:real_perception_pipeline_summary}
\end{figure}

بر اساس توضیحات ارائه‌شده در این بخش، بینایی لایه ادراک سامانه عملی از یک زنجیره چندمرحله‌ای تشکیل شده است که از دریافت تصاویر خام استریو آغاز می‌شود و به تولید خروجی‌های قابل استفاده برای ناوبری ختم می‌گردد. همان‌گونه که در شکل \ref{fig:real_perception_pipeline_summary} نشان داده شده است، تصاویر دو حسگر دوربین استریو ابتدا در رزبری دریافت می‌شوند. این تصاویر در قالب \lr{YUV420} در دسترس قرار گرفته و با استخراج مؤلفه روشنایی، به تصویر تک‌رنگ \lr{mono8} تبدیل می‌شوند. سپس دو جریان تصویری مستقل از طریق \lr{GStreamer} و بستر شبکه محلی به لپ‌تاپ ارسال می‌گردند. در لپ‌تاپ، تصاویر دریافت‌شده در قالب \lr{ROS2} دریافت شده، تصحیح هندسی روی آن‌ها اعمال می‌شود و تصاویر چپ و راست تصحیح‌شده برای مراحل بعدی در دسترس قرار می‌گیرند.

پس از این مرحله، تصاویر تصحیح‌شده مبنای سه خروجی اصلی لایه ادراک قرار می‌گیرند. خروجی نخست، برآورد اختلاف مکانی و عمق است که برای تخمین فاصله هدف و موانع استفاده می‌شود. خروجی دوم، تشخیص موانع است که با استفاده از آشکارساز دیداری و ترکیب آن با نقشه اختلاف مکانی، موقعیت نسبی و در ادامه سرعت نسبی موانع را فراهم می‌کند. خروجی سوم، تشخیص هدف است که در سامانه عملی با استفاده از نشانگر \lr{AprilTag} انجام می‌شود. در مجموع، شکل \ref{fig:real_perception_pipeline_summary} نشان می‌دهد که لایه ادراک در سامانه عملی تنها شامل دریافت تصویر نیست، بلکه یک زنجیره یکپارچه از دریافت، تبدیل قالب، انتقال، تصحیح هندسی و استخراج اطلاعات سطح‌بالا را در بر می‌گیرد که در نهایت داده‌های لازم برای تشکیل مشاهده عامل تصمیم‌گیر را فراهم می‌سازد.


\subsection{پردازش داده حسگرهای فراصوت}
\label{subsec:real_ultrasonic_processing}

در کنار ادراک دیداری مبتنی بر دوربین استریو، از پنج حسگر فراصوت \lr{SRF05} برای پایش موانع نزدیک در اطراف اسکوتر استفاده شد. این حسگرها نقش مکمل دارند و به‌ویژه برای نواحی نزدیک‌بردی به کار می‌روند که واکنش سریع سامانه اهمیت بیشتری دارد.

داده حسگرهای فراصوت توسط آردوینو دریافت می‌شود. برای هر حسگر، آردوینو یک پالس تحریک کوتاه روی پایه تحریک اعمال کرده و مدت زمان پالس برگشتی روی پایه پژواک را اندازه‌گیری می‌کند. در پیاده‌سازی انجام‌شده، پالس تحریک به‌مدت ($10~\mu s$) اعمال شده و فاصله مانع از زمان رفت‌وبرگشت موج صوتی محاسبه می‌شود.

برای جلوگیری از توقف کنترل پایین‌سطح، خواندن حسگرها به‌صورت غیرمسدودکننده انجام شد. به این ترتیب، حسگرها به‌صورت مرحله‌ای و نوبتی خوانده می‌شوند و اندازه‌گیری فراصوت باعث توقف کنترل سرعت موتورها یا دریافت فرمان سریال نمی‌شود. در صورت دریافت‌نشدن پاسخ معتبر از یک حسگر، مقدار آن برابر (-1) در نظر گرفته می‌شود که نشان‌دهنده اندازه‌گیری نامعتبر یا خارج از محدوده قابل اعتماد است.

خروجی این بخش شامل پنج مقدار فاصله متناظر با نواحی نصب حسگرها روی بدنه است. این مقادیر در تشکیل مشاهده عامل و همچنین در لایه ایمنی برای تشخیص وضعیت‌های نزدیک‌برد استفاده می‌شوند. بنابراین، حسگرهای فراصوت اطلاعات هندسی کامل از مانع تولید نمی‌کنند، اما با فراهم‌کردن فاصله نزدیک‌برد در چند جهت مشخص، به افزایش ایمنی سامانه در مجاورت موانع کمک می‌کنند.

\subsection{پردازش داده \lr{IMU}، انکودر و اودومتری نسبی}
\label{subsec:real_imu_encoder_odometry_processing}

در کنار اطلاعات محیطی حاصل از دوربین استریو و حسگرهای فراصوت، لازم است وضعیت حرکتی خود اسکوتر نیز نسبت به نقطه شروع آزمون تخمین زده شود. در پیاده‌سازی عملی، این وضعیت به‌صورت اودومتری نسبی محاسبه شد. منظور از اودومتری نسبی در اینجا، تخمین جابه‌جایی طولی، جابه‌جایی جانبی و زاویه جهت‌گیری اسکوتر نسبت به وضعیت اولیه آزمون است. این اطلاعات برای ثبت رفتار حرکتی سامانه، تحلیل مسیر طی‌شده و تشکیل بخشی از وضعیت ادراکی مورد استفاده قرار می‌گیرند.

در سامانه عملی، سرعت خطی اسکوتر از داده انکودرهای چرخ‌ها و زاویه جهت‌گیری از واحد اندازه‌گیری اینرسی استخراج شد. دلیل استفاده از این ترکیب آن است که در اسکوترهای مبتنی بر رانش لغزشی، تخمین زاویه چرخش فقط از اختلاف سرعت چرخ‌های چپ و راست می‌تواند خطای قابل توجهی داشته باشد. در چنین سامانه‌ای، بخشی از دوران چرخ‌ها صرف لغزش جانبی و سایش روی زمین می‌شود و بنابراین مدل سینماتیکی ایده‌آل نمی‌تواند زاویه واقعی چرخش بدنه را با دقت کافی تخمین بزند. به همین دلیل، در پیاده‌سازی عملی، زاویه \lr{yaw} حسگر اینرسی به‌عنوان منبع اصلی تخمین جهت‌گیری اسکوتر استفاده شد.

داده \lr{yaw} از حسگر \lr{TSS-USB-S} و از طریق ارتباط سریال دریافت می‌شود. در ابتدای هر آزمون، چند نمونه اولیه از \lr{yaw} در حالی که اسکوتر ثابت است دریافت شده و میانگین آن‌ها به‌عنوان مرجع صفر در نظر گرفته می‌شود. سپس مقدار \lr{yaw} در طول آزمون نسبت به این مقدار اولیه بیان می‌گردد. همچنین برای جلوگیری از پرش زاویه هنگام عبور از مرز $-\pi$ و $+\pi$، زاویه دریافتی به‌صورت پیوسته بازگشایی می‌شود. به این ترتیب، جهت‌گیری اسکوتر در طول آزمون به‌صورت یک زاویه نسبی پیوسته نسبت به وضعیت شروع در دسترس قرار می‌گیرد.

سرعت خطی اسکوتر از \lr{RPM} چرخ‌ها محاسبه می‌شود. برای هر سمت اسکوتر، داده‌های چرخ جلو و عقب همان سمت در نظر گرفته می‌شوند. از آنجا که در شرایط لغزش ممکن است یکی از چرخ‌های یک سمت سریع‌تر از حرکت واقعی روی زمین بچرخد، در پیاده‌سازی انجام‌شده مقدار کوچک‌تر بین سرعت چرخ جلو و عقب هر سمت به‌عنوان تخمین محافظه‌کارانه‌تر سرعت همان سمت انتخاب شد. سپس سرعت خطی خام اسکوتر از میانگین سرعت سمت چپ و سمت راست محاسبه می‌شود. در این محاسبات، قطر چرخ برابر $0.185~m$ در نظر گرفته شد.

با داشتن سرعت خطی و زاویه جهت‌گیری، موقعیت نسبی اسکوتر نسبت به نقطه شروع آزمون با انتگرال‌گیری زمانی محاسبه می‌شود. اگر $v$ سرعت خطی اصلاح‌شده و $\theta$ زاویه جهت‌گیری به‌دست‌آمده از \lr{IMU} باشد، تغییرات موقعیت در هر گام زمانی به‌صورت زیر محاسبه می‌شود:

\begin{equation}
	x_{k+1} = x_k + v_k \cos(\theta_k)\Delta t,
	\qquad
	y_{k+1} = y_k + v_k \sin(\theta_k)\Delta t
	\label{eq:real_odom_position_update}
\end{equation}

در این رابطه، $x$ و $y$ مختصات نسبی اسکوتر نسبت به نقطه شروع آزمون هستند و $\Delta t$ بازه زمانی میان دو به‌روزرسانی متوالی است. مقدار مسیر طی‌شده نیز با انتگرال‌گیری از قدرمطلق سرعت خطی در طول زمان محاسبه می‌شود.

در کنار استفاده مستقیم از زاویه \lr{IMU}، تخمین زاویه چرخشی مبتنی بر اختلاف سرعت چرخ‌ها نیز در سامانه محاسبه و ثبت می‌شود. این مقدار به‌عنوان منبع اصلی جهت‌گیری استفاده نمی‌شود، بلکه برای عیب‌یابی، مقایسه و حالت پشتیبان نگه‌داری می‌گردد. در صورتی که داده \lr{IMU} قطع شود یا برای مدت کوتاهی به‌روزرسانی تازه‌ای از آن دریافت نشود، سامانه به‌صورت موقت از تخمین مبتنی بر انکودر برای به‌روزرسانی زاویه استفاده می‌کند.

به دلیل ماهیت رانش لغزشی، سرعت به‌دست‌آمده از انکودرها در هنگام چرخش می‌تواند بیش از جابه‌جایی واقعی اسکوتر باشد. برای کاهش این اثر، یک ضریب اصلاح لغزش به‌صورت تجربی در نظر گرفته شد. این ضریب از مقایسه نرخ چرخش تخمین‌زده‌شده توسط انکودرها با نرخ چرخش اندازه‌گیری‌شده توسط \lr{IMU} به دست می‌آید. در حالتی که چرخش معنادار وجود داشته باشد، نسبت این دو نرخ به‌عنوان شاخصی از میزان تبدیل واقعی دوران چرخ‌ها به دوران بدنه استفاده می‌شود. سپس این ضریب در سرعت خطی خام اعمال می‌شود تا سرعت اصلاح‌شده برای انتگرال‌گیری موقعیت به دست آید. برای جلوگیری از نوسان شدید، ضریب اصلاح لغزش در بازه $[0,1]$ محدود شده و با میانگین‌گیری نمایی هموار می‌شود.

این روش اودومتری، یک تخمین نسبی و آزمایشگاهی از حرکت اسکوتر فراهم می‌کند و جایگزین سامانه مکان‌یابی مطلق نیست. خطای آن می‌تواند در اثر لغزش چرخ‌ها، ناهمواری سطح، خطای انکودر، نویز حسگر اینرسی و خطای تجمعی انتگرال‌گیری در طول زمان افزایش یابد. با این حال، ترکیب سرعت خطی مبتنی بر انکودر با جهت‌گیری مبتنی بر \lr{IMU} باعث شد تخمین مسیر نسبت به استفاده صرف از مدل سینماتیکی چرخ‌ها پایدارتر باشد، به‌ویژه در چرخش‌های رانش لغزشی که اختلاف سرعت چرخ‌ها به‌تنهایی تخمین دقیقی از زاویه واقعی بدنه ارائه نمی‌دهد.

\subsection{تجمیع داده‌های ادراکی و ساخت مشاهده}
\label{subsec:real_observation_construction}

پس از استخراج خروجی‌های لایه ادراک، داده‌های حاصل باید به قالب مشاهده عامل تبدیل شوند. هدف از این مرحله آن است که مدل آموزش‌دیده در محیط شبیه‌سازی، در سامانه عملی نیز ورودی‌ای با همان ساختار، ترتیب و مقیاس عددی دریافت کند. در پیاده‌سازی انجام‌شده، این کار در گره اجرای مدل انجام می‌شود. این گره داده‌های مربوط به هدف، موانع، حسگرهای فراصوت و اودومتری را از موضوع‌های ارتباطی \lr{ROS2} دریافت کرده و مشاهده نهایی را برای ورودی مدل \lr{LSTM-SAC} می‌سازد.

مشاهده عامل از دو بخش اصلی تشکیل شده است: بردار ثابت و توالی موانع. بردار ثابت شامل ۱۵ مؤلفه است و اطلاعات کلی وضعیت هدف، حرکت اسکوتر، حسگرهای فراصوت، فرمان قبلی و وضعیت ایمنی را در خود نگه می‌دارد. ترتیب این مؤلفه‌ها مطابق محیط آموزش به‌صورت زیر حفظ شد:

\begin{equation}
	\begin{aligned}
		\mathbf{o}_{\mathrm{fixed}} =
		[&
		d_g/15,\;
		\beta_g/\pi,\;
		v_{gx}^{\mathrm{rel}}/2,\;
		v_{gy}^{\mathrm{rel}}/2,\;
		v/1,\;
		\omega/1,\;\\
		&
		d_{ufl}/3,\;
		d_{ufr}/3,\;
		d_{uls}/3,\;
		d_{urs}/3,\;
		d_{ub}/3,\;\\
		&
		v_{\mathrm{prev}}/1,\;
		\omega_{\mathrm{prev}}/1,\;
		s_{\mathrm{shield}},\;
		t_{\mathrm{shield}}
		]
	\end{aligned}
	\label{eq:real_fixed_observation_vector}
\end{equation}

در این رابطه، ($d_g$) فاصله هدف، ($\beta_g$) زاویه نسبی هدف، ($ v_{gx}^{rel} $) و ($ v_{gy}^{rel} $) سرعت نسبی هدف، (v) و ($\omega$) سرعت خطی و زاویه‌ای اسکوتر، و ($d_{ufl}$)، ($d_{ufr}$)، ($d_{uls}$)، ($d_{urs}$) و ($d_{ub}$) فاصله‌های حسگرهای فراصوت جلو-چپ، جلو-راست، چپ، راست و عقب هستند. همچنین ($v_{prev}$) و ($\omega_{prev}$) فرمان خطی و زاویه‌ای قبلی عامل را نشان می‌دهند و دو مؤلفه پایانی شدت و مدت فعال‌بودن وضعیت سپر ایمنی را بیان می‌کنند.

بخش دوم مشاهده، ماتریس موانع است. برای هر مانع، پنج ویژگی شامل مختصات نسبی، فاصله و سرعت نسبی در نظر گرفته می‌شود:

\begin{equation}
	\mathbf{o}*{obs,i} =
	[
	x_i/10,;
	y_i/10,;
	d_i/10,;
	v*{xi}^{rel}/2,;
	v_{yi}^{rel}/2
	]
	\label{eq:real_obstacle_observation_vector}
\end{equation}

موانع بر اساس فاصله مرتب می‌شوند و نزدیک‌ترین موانع در ابتدای ماتریس قرار می‌گیرند. تعداد ردیف‌های این ماتریس برابر با بیشینه تعداد موانعی است که مدل با آن آموزش دیده است. در پیاده‌سازی انجام‌شده، این مقدار باید دقیقاً با مقدار استفاده‌شده در آموزش یکسان باشد؛ در غیر این صورت، شکل ورودی مدل با ساختار آموزش‌دیده سازگار نخواهد بود.

برای جلوگیری از گسست ناگهانی در مشاهده موانع، یک حافظه کوتاه‌مدت از موانع در چارچوب جهانی نگه‌داری می‌شود. موانع تشخیص‌داده‌شده ابتدا با استفاده از اودومتری نسبی اسکوتر از چارچوب محلی به چارچوب جهانی منتقل می‌شوند و سپس در مرحله ساخت مشاهده دوباره نسبت به وضعیت فعلی اسکوتر به چارچوب محلی بازگردانده می‌شوند. این حافظه به مدت ($10~s$) اعتبار دارد و باعث می‌شود در صورت حذف کوتاه‌مدت یک مانع از خروجی آشکارساز، مشاهده عامل بلافاصله دچار تغییر ناگهانی نشود.

در نهایت، بردار ثابت و ماتریس موانع به‌صورت یک مشاهده ساختاریافته به مدل داده می‌شوند. تمام مؤلفه‌ها با همان ضرایب نرمال‌سازی محیط آموزش مقیاس‌دهی می‌شوند و مقادیر خروجی در بازه‌های معتبر محدود می‌گردند. این موضوع اهمیت زیادی دارد، زیرا مدل \lr{LSTM-SAC} نسبت به ترتیب و مقیاس ورودی‌ها حساس است و تغییر در این ساختار می‌تواند باعث افت عملکرد در اجرای واقعی شود.

\subsection{مدیریت داده‌های نامعتبر}
\label{subsec:real_invalid_data_handling}

در سامانه عملی، برخلاف محیط شبیه‌سازی، داده‌های حسگری همواره کامل، دقیق و بدون تأخیر نیستند. ممکن است نشانگر هدف برای چند فریم از میدان دید خارج شود، خروجی تشخیص موانع به‌صورت لحظه‌ای حذف گردد، مقدار حسگر فراصوت نامعتبر باشد یا برخی داده‌ها با تأخیر به گره تصمیم‌گیر برسند. از آنجا که عامل \lr{LSTM-SAC} با ساختار مشاهده ثابت آموزش دیده است، پیش از ارسال داده به مدل، این موارد مدیریت می‌شوند تا شکل، ترتیب و محدوده عددی مشاهده حفظ شود.

برای حسگرهای فراصوت، مقدار (-1) به‌عنوان نبود اندازه‌گیری معتبر در نظر گرفته می‌شود. همچنین اگر داده یک حسگر برای مدت مشخصی به‌روزرسانی نشود، مقدار آن قدیمی تلقی می‌گردد. در هر دو حالت، مقدار فاصله آن حسگر برابر با بیشینه برد معتبر جایگزین می‌شود تا داده نامعتبر به‌اشتباه به‌عنوان وجود مانع بسیار نزدیک تفسیر نشود.

در مورد هدف، فاصله و زاویه فقط زمانی معتبر هستند که نشانگر \lr{AprilTag} دیده شود و مقدار عمق آن نیز قابل اعتماد باشد. اگر هدف در یک گام زمانی مشاهده نشود، مشاهده تازه‌ای بر اساس داده هدف ساخته نمی‌شود و آخرین مشاهده معتبر برای ورودی مدل نگه داشته می‌شود. این کار از ایجاد جهش ناگهانی در بردار مشاهده جلوگیری می‌کند، هرچند برای گم‌شدن طولانی‌مدت هدف طراحی نشده است.

برای موانع نیز تازگی داده‌ها بررسی می‌شود. اگر فهرست موانع در بازه زمانی قابل قبول به‌روزرسانی شده باشد، حافظه موانع با داده جدید اصلاح می‌شود؛ در غیر این صورت، وضعیت قبلی تا زمان انقضا حفظ می‌گردد. این روش از حذف ناگهانی موانع در اثر خطای لحظه‌ای آشکارساز جلوگیری می‌کند و هم‌زمان مانع باقی‌ماندن دائمی داده‌های قدیمی در مشاهده عامل می‌شود.

در نهایت، تمام مؤلفه‌های مشاهده پیش از ورود به مدل نرمال‌سازی و در بازه‌های مورد انتظار محدود می‌شوند. بنابراین، مدیریت داده‌های نامعتبر در این بخش با هدف تغییر رفتار عامل انجام نمی‌شود، بلکه برای حفظ قرارداد ورودی مدل در شرایط واقعی به کار می‌رود.

\subsection{اجرای مدل \lr{LSTM-SAC}}
\label{subsec:real_lstm_sac_execution}

پس از تشکیل مشاهده، مدل تصمیم‌گیر آموزش‌دیده روی لپ‌تاپ اجرا می‌شود. اجرای مدل روی لپ‌تاپ انجام شد تا توان پردازشی کافی برای اجرای هم‌زمان شبکه عصبی، پردازش تصویر، تشخیص موانع و ارتباط با سایر گره‌های سامانه فراهم باشد.

در هر گام زمانی، مشاهده ساخته‌شده به مدل داده می‌شود و سیاست با استفاده از تابع پیش‌بینی، کنش مناسب را تولید می‌کند. در اجرای عملی، پیش‌بینی به‌صورت قطعی انجام شد تا در آزمون‌های واقعی رفتار عامل پایدارتر و قابل تکرارتر باشد. خروجی خام مدل یک کنش دوبعدی نرمال‌شده است:

\begin{equation}
	\mathbf{a}*t =
	\begin{bmatrix}
		a_v \
		a*\omega
	\end{bmatrix},
	\qquad
	a_v,a_\omega \in [-1,1]
	\label{eq:real_policy_action}
\end{equation}

که در آن ($a_v$) مؤلفه مربوط به سرعت خطی و ($a_\omega$) مؤلفه مربوط به سرعت زاویه‌ای است. این مؤلفه‌ها با همان منطق تعریف‌شده در فصل چهارم به سرعت خطی و زاویه‌ای مطلوب تبدیل می‌شوند:

\begin{equation}
	v_{des} = a_v v_{max},
	\qquad
	\omega_{des} = a_\omega \omega_{max}
	\label{eq:real_action_scaling}
\end{equation}

در پیاده‌سازی انجام‌شده، مقدارهای ($v_{max}=1.0~m/s$) و ($\omega_{max}=1.0~rad/s$) استفاده شدند تا مقیاس خروجی مدل با محیط آموزش هم‌خوان باقی بماند.

برای جلوگیری از تغییرات ناگهانی فرمان، خروجی مدل پیش از انتشار هموارسازی شد. این هموارسازی به‌صورت میانگین‌گیری نمایی میان فرمان قبلی و فرمان جدید انجام می‌شود:

\begin{equation}
	v_t =
	(1-\alpha)v_{t-1}
	+
	\alpha v_{des}, \quad
	\omega_t =
	(1-\alpha)\omega_{t-1}
	+
	\alpha \omega_{des}
	\label{eq:real_w_and_v_smoothing}
\end{equation}

که در آن ($\alpha=0.70$) در نظر گرفته شد. با این کار، فرمان نهایی همچنان به تصمیم جدید عامل وابسته است، اما جهش‌های لحظه‌ای خروجی کاهش می‌یابد و فرمان حرکتی برای اجرا روی سامانه واقعی نرم‌تر می‌شود.

خروجی نهایی این گره به‌صورت پیام \lr{Twist} روی موضوع ارتباطی سرعت ربات منتشر می‌شود. این پیام شامل سرعت خطی ($v_t$) و سرعت زاویه‌ای ($\omega_t$) است و در بخش بعدی به فرمان قابل اجرا برای چهار موتور اسکوتر تبدیل می‌شود.

\subsection{تبدیل خروجی مدل به فرمان حرکتی}
\label{subsec:real_action_to_motion_command}

خروجی گره اجرای مدل به‌صورت پیام \lr{Twist} منتشر می‌شود و شامل سرعت خطی (v) و سرعت زاویه‌ای ($\omega$) است. از آنجا که اسکوتر دارای چهار موتور مستقل است، این فرمان سطح‌بالا باید پیش از اجرا به سرعت هدف موتورهای چهارگانه تبدیل شود.

این تبدیل در گره واسط سخت‌افزاری لپ‌تاپ انجام می‌شود. ابتدا سرعت معادل سمت چپ و راست اسکوتر بر اساس مدل رانش تفاضلی محاسبه می‌گردد:

\begin{equation}
	v_L = v - \omega \frac{L_{eff}}{2},
	\qquad
	v_R = v + \omega \frac{L_{eff}}{2}
	\label{eq:real_left_and_right_side_velocity}
\end{equation}

در این رابطه، ($L_{eff}$) فاصله مؤثر میان دو سمت حرکتی است که در پیاده‌سازی عملی برابر ($0.66108~m$) در نظر گرفته شد. سپس با استفاده از قطر چرخ ($D=0.185~m$)، سرعت خطی هر سمت به \lr{RPM} تبدیل می‌شود:

\begin{equation}
	RPM = \frac{v_{wheel}}{\pi D} \times 60
	\label{eq:real_wheel_speed_to_rpm}
\end{equation}

بر این اساس، چرخ‌های جلو و عقب هر سمت مقدار هدف یکسانی دریافت می‌کنند. با توجه به جهت نصب موتورهای سمت راست، علامت فرمان این موتورها در قالب نهایی معکوس می‌شود و بردار فرمان چهار موتور به‌صورت زیر ساخته می‌شود:

\begin{equation}
	\mathbf{u}_{rpm} =
	[
	RPM_L,;
	-RPM_R,;
	RPM_L,;
	-RPM_R
	]
	\label{eq:real_four_motor_rpm_command}
\end{equation}

ترتیب مؤلفه‌ها به‌ترتیب متناظر با موتور جلو-چپ، جلو-راست، عقب-چپ و عقب-راست است. برای جلوگیری از اعمال فرمان بیش از حد، مقدار \lr{RPM} هر موتور در محدوده مجاز محدود می‌شود؛ در پیاده‌سازی عملی، بیشینه مقدار مجاز برابر ($150~RPM$) در نظر گرفته شد.

پس از تشکیل فرمان چهارموتوره، لپ‌تاپ آن را از طریق \lr{UDP} به رزبری پای ارسال می‌کند. رزبری پای فرمان دریافتی را به قالب سریال قابل فهم برای آردوینو تبدیل کرده و از طریق \lr{USB/Serial} به آن منتقل می‌کند. در نهایت، آردوینو با استفاده از بازخورد انکودرها، کنترل سرعت پایین‌سطح هر موتور را انجام داده و سیگنال \lr{PWM} مناسب را برای راه‌اندازهای موتور تولید می‌کند. برای جلوگیری از باقی‌ماندن آخرین فرمان در شرایط خطا، در مسیر ارتباطی فرمان، سازوکار توقف خودکار نیز در نظر گرفته شد. در صورت قطع ارتباط یا دریافت‌نشدن فرمان تازه، فرمان صفر برای موتورهای چهارگانه ارسال می‌شود.
\section{پیاده‌سازی کنترل پایین‌سطح و ایمنی در سامانه عملی}
\label{sec:real_low_level_control_and_safety}

پس از تولید فرمان حرکتی توسط مدل تصمیم‌گیر، لازم است این فرمان به‌صورت قابل اجرا و ایمن به موتورهای اسکوتر اعمال شود. در سامانه عملی، فرمان عامل مستقیماً به راه‌اندازهای موتور ارسال نمی‌شود، بلکه ابتدا به سرعت هدف چهار موتور تبدیل شده و سپس در آردوینو، با استفاده از بازخورد انکودرها، کنترل سرعت پایین‌سطح انجام می‌گیرد. علاوه بر این، برای کاهش خطر برخورد و جلوگیری از ادامه حرکت در شرایط خطا، منطق ایمنی و سازوکار توقف خودکار نیز در مسیر فرمان در نظر گرفته شد.

\subsection{دریافت فرمان و کنترل سرعت موتورهای چهارگانه}
\label{subsec:real_motor_speed_control}

فرمان چهارموتوره پس از تولید در گره واسط سخت‌افزاری، از رزبری پای به آردوینو ارسال می‌شود. این فرمان شامل سرعت هدف هر چهار موتور بر حسب \lr{RPM} است و ترتیب مقادیر آن متناظر با موتورهای جلو-چپ، جلو-راست، عقب-چپ و عقب-راست در نظر گرفته می‌شود. آردوینو پس از دریافت فرمان، مقدار هدف هر موتور را استخراج کرده و آن را به‌عنوان ورودی حلقه کنترل سرعت استفاده می‌کند.

هر موتور دارای انکودر مستقل است و آردوینو با شمارش پالس‌های انکودر، سرعت واقعی موتور را تخمین می‌زند. سپس برای هر موتور، خطای سرعت از اختلاف میان سرعت هدف و سرعت اندازه‌گیری‌شده محاسبه می‌شود:

\begin{equation}
	e_i = RPM^{ref}_i - RPM^{meas}_i
	\label{eq:real_motor_rpm_error}
\end{equation}

در این رابطه، ($RPM^{ref}_i$) سرعت هدف موتور (i)، ($RPM^{meas}_i$) سرعت اندازه‌گیری‌شده و ($e_i$) خطای سرعت است. بر اساس این خطا، سیگنال \lr{PWM} و جهت چرخش مناسب برای راه‌انداز موتور تولید می‌شود. بنابراین، آردوینو صرفاً یک فرمان باز به موتورها اعمال نمی‌کند، بلکه با استفاده از بازخورد انکودر، سرعت هر موتور را به مقدار هدف نزدیک می‌سازد.

کنترل مستقل چهار موتور برای اسکوتر چهارچرخ محرک اهمیت دارد، زیرا تفاوت در بار چرخ‌ها، اصطکاک سطح، پاسخ موتورها و جهت نصب آن‌ها می‌تواند باعث اختلاف میان سرعت واقعی چرخ‌ها شود. استفاده از حلقه کنترل سرعت مبتنی بر انکودر، این اختلاف‌ها را تا حدی کاهش داده و اجرای فرمان حرکتی را پایدارتر می‌کند.

\subsection{اعمال منطق ایمنی در مسیر فرمان}
\label{subsec:real_safety_logic_execution}

در سامانه عملی، فرمان تولیدشده توسط عامل تصمیم‌گیر پیش از اجرا از نظر ایمنی بررسی می‌شود. هدف این لایه، تغییر سیاست یادگیری‌شده نیست، بلکه جلوگیری از اعمال فرمان‌هایی است که در شرایط واقعی می‌توانند باعث نزدیک‌شدن بیش از حد اسکوتر به مانع یا برخورد شوند. منطق ایمنی مورد استفاده، همان ساختار سپر ایمنی تعریف‌شده در فصل چهارم را دنبال می‌کند؛ با این تفاوت که در اجرای عملی، اثر هر سطح خطر به‌صورت مستقیم روی فرمان سرعت خطی و زاویه‌ای اعمال می‌شود.

در این منطق، داده حسگرهای فراصوت برای تشخیص موانع بسیار نزدیک و خروجی تشخیص دیداری برای بررسی موانع روبه‌روی اسکوتر استفاده می‌شود. اگر هیچ مانع نزدیکی در محدوده خطر وجود نداشته باشد، فرمان عامل بدون تغییر اجرا می‌شود. اما در صورت ورود به نواحی احتیاط، خطر یا بحرانی، فرمان نهایی مطابق الگوریتم \ref{alg:real_safety_shield_execution} اصلاح می‌گردد.
\begin{algorithm}[htb]
	\footnotesize
	\caption{اعمال سپر ایمنی فعال روی فرمان حرکتی عامل}
	\label{alg:real_safety_shield_execution}
	\begin{algorithmic}[1]
		\REQUIRE فرمان عامل $\mathbf{u}*{policy}=(v*{policy},\omega_{policy})$، فاصله حسگرهای فراصوت، فاصله نزدیک‌ترین مانع دیداری روبه‌رو
		\ENSURE فرمان نهایی قابل اجرا $\mathbf{u}*{final}=(v*{final},\omega_{final})$
		
		\STATE $v_{final}\gets v_{policy}$ ، $\omega_{final}\gets \omega_{policy}$
		\STATE $d_{front}\gets \min(d_{FL},d_{FR})$
		\STATE $d_{side}\gets \min(d_{LS},d_{RS})$
		\STATE $d_{all}\gets \min(d_{front},d_{side}, d_{back})$
		\STATE $d_{left}\gets \min(d_{FL},d_{LS})$ ، $d_{right}\gets \min(d_{FR},d_{RS})$
		\IF{$d_{all}<0.25~m$ یا $d^{obs}_{front}<0.9~m$}
		\IF{$d_{back}>0.5~m$}
		\STATE $v_{final}\gets -0.15~m/s$
		\ELSE
		\STATE $v_{final}\gets 0$
		\ENDIF
		\IF{$d_{left}\geq d_{right}$}
		\STATE $\omega_{final}\gets +0.8~rad/s$
		\ELSE
		\STATE $\omega_{final}\gets -0.8~rad/s$
		\ENDIF
		\RETURN $\mathbf{u}_{final}$
		\ENDIF
		\IF{$d_{front}<0.50~m$ یا $d^{obs}_{front}<1.5~m$}
		\STATE $v_{final}\gets \min(v_{final},0.10~m/s)$
		\IF{مانع نزدیک‌تر در سمت چپ باشد}
		\STATE $\omega_{final}\gets \omega_{final}-0.5~rad/s$
		\ELSE
		\STATE $\omega_{final}\gets \omega_{final}+0.5~rad/s$
		\ENDIF
		\STATE $\omega_{final}\gets \mathrm{clip}(\omega_{final},-1,1)$
		\RETURN $\mathbf{u}_{final}$
		\ENDIF
		\IF{$d_{front}<0.80~m$ یا $d^{obs}_{front}<2.5~m$ یا $d_{side}<0.50~m$}
		\STATE $v_{final}\gets \min(v_{final},0.40~m/s)$
		\RETURN $\mathbf{u}_{final}$
		\ENDIF
		\RETURN $\mathbf{u}_{policy}$ 
	\end{algorithmic}
	
\end{algorithm}



در سطح احتیاط، هدف توقف کامل اسکوتر نیست، بلکه کاهش سرعت پیشروی است تا سامانه فرصت بیشتری برای اصلاح مسیر داشته باشد. به همین دلیل، سرعت خطی مثبت به حداکثر $0.40~m/s$ محدود می‌شود و فرمان چرخشی عامل تغییر نمی‌کند. در سطح خطر، سرعت پیشروی شدیدتر محدود شده و حداکثر مقدار آن $0.10~m/s$ در نظر گرفته می‌شود. همچنین یک بایاس چرخشی با مقدار $0.5~rad/s$ به فرمان اضافه می‌شود تا اسکوتر به سمت فضای بازتر حرکت کند.

در سطح بحرانی، هدف اصلی خروج سریع از وضعیت برخورد احتمالی است. در این حالت، اگر فضای عقب آزاد باشد، فرمان عقب‌روی آرام با مقدار $-0.15~m/s$ اعمال می‌شود و هم‌زمان اسکوتر با سرعت زاویه‌ای $0.8~rad/s$ به سمت فضای بازتر می‌چرخد. اگر حسگر عقب نشان دهد که فضای کافی برای عقب‌روی وجود ندارد، مؤلفه عقب‌روی حذف شده و اسکوتر فقط درجا به سمت فضای بازتر می‌چرخد. این شرط برای جلوگیری از برخورد عقب اسکوتر با مانع پشت سر در نظر گرفته شده است.

به‌صورت کلی، فرمان نهایی قابل اجرا را می‌توان به‌صورت زیر بیان کرد:

\begin{equation}
	\mathbf{u}_{final} =
	\begin{cases}
		\mathbf{u}_{policy}, & \text{وضعیت عادی} \\
		\mathbf{u}_{limited}, & \text{وضعیت احتیاط} \\
		\mathbf{u}_{safe}, & \text{وضعیت خطر یا بحرانی}
	\end{cases}
	\label{eq:real_safe_command_selection}
\end{equation}

که در آن $\mathbf{u}_{policy}$ فرمان تولیدشده توسط عامل، $\mathbf{u}_{limited}$ نسخه محدودشده فرمان عامل و $\mathbf{u}_{safe}$ فرمان اصلاح‌شده برای خروج از وضعیت خطر است. بنابراین، سپر ایمنی به‌عنوان یک ناظر روی مسیر فرمان عمل می‌کند؛ عامل همچنان تصمیم اصلی را تولید می‌کند، اما در شرایطی که اجرای مستقیم آن با ایمنی سخت‌افزار واقعی سازگار نباشد، فرمان نهایی به مقدار محافظه‌کارانه‌تر تغییر داده می‌شود.

علاوه بر این، برای جلوگیری از گیرکردن اسکوتر در وضعیت توقف، یک منطق خروج از بن‌بست نیز در نظر گرفته شد. اگر فرمان نهایی برای چند گام متوالی تقریباً صفر باقی بماند، سامانه یک چرخش اجباری به سمت فضای بازتر اعمال می‌کند و در صورت خالی‌بودن عقب، مقدار کمی عقب‌روی نیز به آن اضافه می‌شود. این رفتار باعث می‌شود اسکوتر در نزدیکی موانع فقط متوقف نماند و بتواند از وضعیت‌های بحرانی یا بن‌بست خارج شود.

\section{آماده‌سازی و اجرای آزمون‌های عملی}
\label{sec:real_test_preparation_and_execution}

پس از پیاده‌سازی زنجیره ادراک، تصمیم‌گیری، کنترل پایین‌سطح و ایمنی، آزمون‌های عملی برای بررسی عملکرد یکپارچه سامانه طراحی شدند. هدف از این بخش، ارائه نتایج نهایی نیست، بلکه توضیح نحوه آماده‌سازی سناریوها و ثبت داده‌هایی است که در فصل ششم برای تحلیل عملکرد استفاده می‌شوند.

\subsection{سناریوهای آزمون و نحوه اجرای آن‌ها}
\label{subsec:real_test_scenarios_execution}

سناریوهای آزمون عملی به‌گونه‌ای طراحی شدند که اجزای اصلی سامانه، شامل تشخیص هدف، تشخیص موانع، تشکیل مشاهده، اجرای مدل، ارسال فرمان و کنترل موتورهای چهارگانه، در قالب یک زنجیره کامل فعال باشند. در ساده‌ترین سناریو، هدف مجهز به نشانگرهای \lr{AprilTag} در فاصله‌ای مشخص روبه‌روی اسکوتر قرار گرفت و سامانه در محیط بدون مانع اجرا شد. این آزمون برای بررسی صحت تشخیص هدف، علامت زاویه نسبی، جهت حرکت چرخ‌ها و پیوستگی فرمان حرکتی استفاده شد.

در گروه بعدی آزمون‌ها، موقعیت هدف نسبت به اسکوتر تغییر داده شد. هدف می‌توانست با زاویه جانبی نسبت به محور روبه‌روی اسکوتر قرار گیرد یا در طول آزمون جابه‌جا شود. این سناریو برای بررسی واکنش سامانه نسبت به تغییرات زاویه و فاصله هدف و همچنین پایداری خروجی مدل در شرایط واقعی به کار رفت.

برای بررسی نقش ادراک موانع، سناریوهایی با حضور مانع در مسیر حرکت اسکوتر تعریف شد. در این آزمون‌ها، خروجی آشکارساز دیداری، نقشه اختلاف مکانی و فیلتر کالمن برای تخمین موقعیت و سرعت نسبی مانع مورد استفاده قرار گرفتند. هدف از این گروه آزمون‌ها، ثبت رفتار سامانه در شرایطی بود که عامل علاوه بر دنبال‌کردن هدف، باید وجود مانع را نیز در مشاهده خود دریافت کند.

در سناریوهای نزدیک‌برد، موانع در فاصله‌های کم و در نواحی پوشش حسگرهای فراصوت قرار داده شدند. این آزمون‌ها برای بررسی نقش حسگرهای فراصوت و سپر ایمنی طراحی شدند؛ به‌ویژه در شرایطی که مانع در نواحی جلو، جانبی یا عقب اسکوتر قرار دارد و لازم است فرمان عامل محدود یا اصلاح شود.

در تمام آزمون‌ها، اجرای سامانه ابتدا با سرعت پایین و در محیط کنترل‌شده آغاز شد. پس از بررسی جداگانه عملکرد هر بخش، زنجیره کامل شامل دریافت تصویر، تشخیص هدف و موانع، تشکیل مشاهده، اجرای مدل، ارسال فرمان و کنترل پایین‌سطح فعال گردید. تحلیل کمی و کیفی نتایج این آزمون‌ها در فصل ششم ارائه می‌شود.

\subsection{ثبت داده‌ها برای تحلیل عملکرد}
\label{subsec:real_test_data_logging}

برای تحلیل رفتار سامانه در آزمون‌های عملی، داده‌های اصلی هر اجرا ثبت شدند. این داده‌ها شامل وضعیت هدف، اطلاعات موانع، مقادیر حسگرهای فراصوت، اودومتری نسبی اسکوتر، فرمان خروجی مدل، فرمان نهایی ارسال‌شده به موتورها و وضعیت لایه ایمنی هستند. ثبت این داده‌ها امکان بررسی دقیق‌تر رفتار سامانه پس از آزمون را فراهم می‌کند.

از لایه ادراک، فاصله و زاویه هدف، وضعیت قابل مشاهده بودن نشانگر \lr{AprilTag}، اطلاعات موانع شامل مختصات نسبی، عمق و سرعت نسبی، و همچنین مقادیر حسگرهای فراصوت ذخیره می‌شوند. این داده‌ها برای بررسی کیفیت مشاهده عامل و تحلیل خطاهای احتمالی در تشخیص هدف یا موانع مورد استفاده قرار می‌گیرند.

از لایه تصمیم‌گیری و کنترل، فرمان خطی و زاویه‌ای تولیدشده توسط مدل، فرمان هموارشده، فرمان نهایی پس از اعمال محدودیت‌های ایمنی، سرعت واقعی چرخ‌ها، مقادیر \lr{RPM} و اودومتری نسبی اسکوتر ثبت می‌شوند. این اطلاعات امکان مقایسه میان فرمان مطلوب و پاسخ واقعی سخت‌افزار را فراهم می‌کند.

در برخی آزمون‌ها، خروجی‌های تصویری و ویدئویی نیز برای تحلیل کیفی ذخیره شدند. این خروجی‌ها شامل تصویر تصحیح‌شده دوربین، نتیجه تشخیص هدف، کادرهای تشخیص موانع و نقشه اختلاف مکانی هستند. همچنین لام به ذکر است، داده‌های ثبت‌شده در این بخش مبنای تحلیل‌های فصل ششم قرار می‌گیرند.



